% notes/artefacts/market-making-lob/derivations.tex
% Full derivations from the legacy notes, kept verbatim when the blocks were migrated on 2026-10-14
% (format pilot, notes/migration.md). Only figure paths changed. The compendium entries state each
% result with its locator and an idea of the proof; the complete working lives here.
% Source of each block as found by the pilot: Cartea, Jaimungal and Penalva (crypto-cartea),
% chapter "A Primer on the Microstructure of Financial Markets", pp. 20-30, except the Kyle
% equilibrium values, which that text misprints (see the entry mod:kyle-single-auction), and the
% Avellaneda-Stoikov blocks (s6-stoikov). Treat everything here as unchecked own working.
% Compile inside notes/artefacts/market-making-lob/ with  latexmk -pdf derivations.tex
\PassOptionsToPackage{table}{xcolor}
\documentclass[a4paper,oneside,openany]{book}
% Your own packages, macros and environments. The learning guide's entry environments,
% \citep and \qfsource are in house.tex, which compendium.tex loads after this file.

% -----------------------------
% Document Class & Encoding
% -----------------------------
\usepackage[utf8]{inputenc}  % Input encoding
\usepackage[T1]{fontenc}     % Font encoding

% -----------------------------
% Page Layout & Geometry
% -----------------------------
\usepackage[margin=2cm,foot=1cm]{geometry}

% -----------------------------
% Header & Footer Management
% -----------------------------
\usepackage{fancyhdr}
\pagestyle{fancy}

% -----------------------------
% Text Formatting & Spacing
% -----------------------------
\usepackage{ragged2e}
\usepackage[nodisplayskipstretch]{setspace}
\setstretch{1}  % Line spacing

% -----------------------------
% Listings for Code
% -----------------------------
\usepackage{listings}
\usepackage{xcolor}

% -----------------------------
% Tables, Figures & Graphics
% -----------------------------
\usepackage{tabularx}
\usepackage{multirow}
\usepackage{pdflscape}
\usepackage{graphicx}

% -----------------------------
% Algorithms
% -----------------------------
\usepackage{algorithm}       % For floating algorithm environment
\usepackage{algpseudocode}   % For pseudo-code inside algorithms

% -----------------------------
% Math Packages
% -----------------------------
\usepackage{amsmath,amssymb,amsthm,mathtools,bm}
\usepackage{mathrsfs}        % Script fonts
\usepackage{tikz-cd}         % Commutative diagrams
\usepackage{rotating}        % Rotated objects
\usepackage{float}           % Float placement
\usetikzlibrary{calc, arrows, positioning, arrows.meta}

% -----------------------------
% TikZ and its Libraries
% -----------------------------
\usepackage{tikz}
\usetikzlibrary{positioning} % Added to support "below=of" syntax
\usepackage{forest}          % For advanced tree diagrams
\usepackage{tikz-qtree}      % For vertical tree diagrams

% -----------------------------
% Miscellaneous Packages
% -----------------------------
\usepackage{xparse}          % For defining robust commands
\usepackage{enumitem}        % For customizing list layouts
\usepackage{makecell}        % For enhanced table cells
\usepackage{comment}		 % For commenting out sections

%%%%%%%%%%%%%%%%%%%%%%%%%%%%%%%%%%%%%%%%%%%%%%%%%%%%%%%%%%%%%%%%%%%%%%%%%%%%%%%
% Custom Math Commands and Shortcuts
%%%%%%%%%%%%%%%%%%%%%%%%%%%%%%%%%%%%%%%%%%%%%%%%%%%%%%%%%%%%%%%%%%%%%%%%%%%%%%%

% Matrix and Vector Formatting
\newcommand{\matr}[1]{\mathbf{#1}}      % For undergraduate algebra (alternative: \bm{#1})
\newcommand{\vect}[1]{\bm{#1}}

% Standard Number Sets
\newcommand{\N}{\mathbb{N}}
\newcommand{\R}{\mathbb{R}}
\newcommand{\Q}{\mathbb{Q}}
\newcommand{\C}{\mathbb{C}}
\newcommand{\Z}{\mathbb{Z}}

% Bold variables (customize as needed)
\newcommand{\x}{\mathbf{x}}
\newcommand{\F}{\mathbf{F}}
\newcommand{\f}{\mathbf{f}}
\newcommand{\y}{\mathbf{y}}
\renewcommand{\a}{\mathbf{a}}
\renewcommand{\b}{\mathbf{b}}
\renewcommand{\c}{\mathbf{c}}
\newcommand{\h}{\mathbf{h}}
\newcommand{\g}{\mathbf{g}}
\newcommand{\z}{\mathbf{z}}
\newcommand{\ze}{\mathbf{0}}

% Norms, Absolute Values, and Bracketing
\newcommand{\norm}[1]{\left\lVert#1\right\rVert}
\newcommand{\abs}[1]{\left\lvert#1\right\rvert}
\newcommand{\brk}[1]{\left[#1\right]}
\newcommand{\brc}[1]{\left\{#1\right\}}
\newcommand{\paren}[1]{\left(#1\right)}
\newcommand{\normop}[1]{\left\lVert#1\right\rVert_\text{op}}

% Miscellaneous Math Shortcuts
\newcommand{\LL}{\mathcal{L}}
\newcommand{\uni}{\overset{\text{uni}}{\to}}
\DeclareMathOperator{\diam}{diam}
\newcommand{\Prr}[1]{\text{Pr}\left(#1\right)}
\newcommand{\ceil}[1]{\lceil #1 \rceil}
\newcommand{\floor}[1]{\lfloor #1 \rfloor}
\newcommand\disteq{\stackrel{\textnormal{dist}}{=}}

% Evaluation command with options
\NewDocumentCommand{\evalat}{sO{\big}mm}{%
  \IfBooleanTF{#1}
    {\mleft. #3 \mright|_{#4}}
    {#3#2|_{#4}}%
}

%%%%%%%%%%%%%%%%%%%%%%%%%%%%%%%%%%%%%%%%%%%%%%%%%%%%%%%%%%%%%%%%%%%%%%%%%%%%%%%
% Color Definitions
%%%%%%%%%%%%%%%%%%%%%%%%%%%%%%%%%%%%%%%%%%%%%%%%%%%%%%%%%%%%%%%%%%%%%%%%%%%%%%%

\definecolor{dartmouthgreen}{rgb}{0.05, 0.5, 0.06}
\definecolor{egyptianblue}{rgb}{0.06, 0.2, 0.65}
\definecolor{dukeblue}{rgb}{0.0, 0.0, 0.61}
\definecolor{jazzberryjam}{rgb}{0.65, 0.04, 0.37}
\definecolor{magenta}{HTML}{EC008C}
\definecolor{darkmagenta}{rgb}{0.55, 0.0, 0.55}
\definecolor{deeppink}{rgb}{1.0, 0.08, 0.58}
\definecolor{codegreen}{rgb}{0,0.6,0}
\definecolor{codegray}{rgb}{0.5,0.5,0.5}
\definecolor{codepurple}{rgb}{0.58,0,0.82}
\definecolor{backcolour}{rgb}{0.95,0.95,0.92}

%%%%%%%%%%%%%%%%%%%%%%%%%%%%%%%%%%%%%%%%%%%%%%%%%%%%%%%%%%%%%%%%%%%%%%%%%%%%%%%
% Custom Shortcuts and Emphasis Commands
%%%%%%%%%%%%%%%%%%%%%%%%%%%%%%%%%%%%%%%%%%%%%%%%%%%%%%%%%%%%%%%%%%%%%%%%%%%%%%%

\newcommand{\isomorp}{\xrightarrow{\sim}}
\newcommand{\hlt}[1]{\textit{{\color{blue}#1}}}
\newcommand{\simpt}[1]{\textit{{\color{deeppink}#1}}}
\newcommand{\impt}[1]{\textit{{\color{red}#1}}}
\newcommand{\gcds}[1]{\textnormal{gcd}#1}
\newcommand{\mins}[1]{\textnormal{min}#1}
\newcommand{\maxs}[1]{\textnormal{max}#1}
\newcommand{\lcms}[1]{\textnormal{lcm}#1}
\newcommand{\degs}[1]{\textnormal{deg}#1}
\newcommand{\exps}[1]{\textnormal{exp}#1}
\newcommand{\tors}[1]{\textnormal{Tor}#1}
\newcommand{\homs}[1]{\textnormal{Hom}#1}
\newcommand{\anns}[1]{\textnormal{Ann}#1}
\newcommand{\defeq}{\vcentcolon=}
\DeclareMathOperator{\sign}{sign}

%%%%%%%%%%%%%%%%%%%%%%%%%%%%%%%%%%%%%%%%%%%%%%%%%%%%%%%%%%%%%%%%%%%%%%%%%%%%%%%
% Theorem and Definition Environments
%%%%%%%%%%%%%%%%%%%%%%%%%%%%%%%%%%%%%%%%%%%%%%%%%%%%%%%%%%%%%%%%%%%%%%%%%%%%%%%
\theoremstyle{plain}
\newtheorem{theorem}{Theorem}[section]
\newtheorem{corollary}[theorem]{Corollary}
\newtheorem{lemma}[theorem]{Lemma}
\newtheorem{proposition}[theorem]{Proposition}
\newtheorem{qnbank}[theorem]{Question Bank}
\let\oldqnbank\qnbank
\renewcommand{\qnbank}{\oldqnbank\normalfont}
\newtheorem{process}[theorem]{Process}
\let\oldprocess\process
\renewcommand{\process}{\oldprocess\normalfont}
\newtheorem{method}[theorem]{Method}
\let\oldmethod\method
\renewcommand{\method}{\oldmethod\normalfont}
\newtheorem{definition}[theorem]{Definition}
\let\olddefinition\definition
\renewcommand{\definition}{\olddefinition\normalfont}
\newtheorem{example}[theorem]{Example}
\let\oldexample\example
\renewcommand{\example}{\oldexample\normalfont}
\newtheorem{remark}[theorem]{Remark}
\let\oldremark\remark
\renewcommand{\remark}{\oldremark\normalfont}
\newtheorem{notation}[theorem]{Notation}
\let\oldnotation\notation
\renewcommand{\notation}{\oldnotation\normalfont}
\newtheorem{principle}[theorem]{Principle}
\let\oldprinciple\principle
\renewcommand{\principle}{\oldprinciple\normalfont}
\newtheorem{mechanism}[theorem]{Mechanism}
\let\oldmechanism\mechanism
\renewcommand{\mechanism}{\oldmechanism\normalfont}
\newtheorem{taxonomy}[theorem]{Taxonomy}
\let\oldtaxonomy\taxonomy
\renewcommand{\taxonomy}{\oldtaxonomy\normalfont}
\newtheorem{model}[theorem]{Model}
\let\oldmodel\model
\renewcommand{\model}{\oldmodel\normalfont}
\newtheorem{assumption}[theorem]{Assumption}
\let\oldassumption\assumption
\renewcommand{\assumption}{\oldassumption\normalfont}
\newtheorem{casestudy}[theorem]{Case Study}
\let\oldcasestudy\casestudy
\renewcommand{\casestudy}{\oldcasestudy\normalfont}
\newtheorem{empiricalfinding}[theorem]{Empirical Finding}
\let\oldempiricalfinding\empiricalfinding
\renewcommand{\empiricalfinding}{\oldempiricalfinding\normalfont}
\newtheorem{practitionernote}[theorem]{Practitioner Note}
\let\oldpractitionernote\practitionernote
\renewcommand{\practitionernote}{\oldpractitionernote\normalfont}
\newtheorem{caveat}[theorem]{Caveat}
\let\oldcaveat\caveat
\renewcommand{\caveat}{\oldcaveat\normalfont}
\newtheorem{marketrule}[theorem]{Rule}
\let\oldmarketrule\marketrule
\renewcommand{\marketrule}{\oldmarketrule\normalfont}
\newtheorem{takeaway}[theorem]{Takeaway}
\let\oldtakeaway\takeaway
\renewcommand{\takeaway}{\oldtakeaway\normalfont}

%%%%%%%%%%%%%%%%%%%%%%%%%%%%%%%%%%%%%%%%%%%%%%%%%%%%%%%%%%%%%%%%%%%%%%%%%%%%%%%
% Breakable Algorithm Environment (Algorithms spanning multiple pages)
%%%%%%%%%%%%%%%%%%%%%%%%%%%%%%%%%%%%%%%%%%%%%%%%%%%%%%%%%%%%%%%%%%%%%%%%%%%%%%%

\makeatletter
\newenvironment{breakablealgorithm}
 {%
  \begin{center}
    \refstepcounter{algorithm} % Increment algorithm counter
    \hrule height.8pt depth0pt \kern2pt
    \renewcommand{\caption}[2][\relax]{%
      {\raggedright\textbf{Algorithm \thealgorithm:} ##2\par}%
      \ifx\relax##1\relax
        \addcontentsline{loa}{algorithm}{\protect\numberline{\thealgorithm}##2}%
      \else
        \addcontentsline{loa}{algorithm}{\protect\numberline{\thealgorithm}##1}%
      \fi
      \kern2pt\hrule\kern2pt
    }%
 }{%
    \kern2pt\hrule
  \end{center}
 }
\makeatother

%%%%%%%%%%%%%%%%%%%%%%%%%%%%%%%%%%%%%%%%%%%%%%%%%%%%%%%%%%%%%%%%%%%%%%%%%%%%%%%
% Code Listings Style
%%%%%%%%%%%%%%%%%%%%%%%%%%%%%%%%%%%%%%%%%%%%%%%%%%%%%%%%%%%%%%%%%%%%%%%%%%%%%%%

\lstdefinestyle{codestyle}{
  commentstyle=\color{codegreen},
  keywordstyle=\color{magenta},
  numberstyle=\tiny\color{codegray},
  stringstyle=\color{codepurple},
  basicstyle=\ttfamily\normalsize,
  breakatwhitespace=false,
  breaklines=true,
  captionpos=b,
  keepspaces=true,
  numbers=left,
  numbersep=10pt,
  frame=single,
  showspaces=false,
  showstringspaces=false,
  showtabs=false,
  tabsize=2
}
\lstset{style=codestyle}

% -----------------------------
% Hyperlinks (loaded last: loaded before float and algorithm it duplicated figure anchors)
% -----------------------------
\usepackage{hyperref}
\hypersetup{
  colorlinks,
  linkcolor={blue!80!black},
  citecolor={green!80!black},
  urlcolor={red!80!black}
}

\usepackage[T1]{fontenc}
\usepackage{amsmath,amssymb,booktabs,graphicx,hyperref}
\hypersetup{colorlinks}
\setlength{\emergencystretch}{3em}   % lets a line holding a long key break
\newcounter{qfentry}[chapter]
\renewcommand{\theqfentry}{\thechapter.\arabic{qfentry}}
% \qftype{type}{Name} declares the environment qf<type>; its title wraps like a paragraph.
\newcommand{\qftype}[2]{\newenvironment{qf#1}[1][]{\par\addvspace{\medskipamount}%
  \refstepcounter{qfentry}{\raggedright\noindent\bfseries #2 \theqfentry
  \if\relax\detokenize{##1}\relax\else\quad ##1\fi\par}\nopagebreak}%
  {\par\addvspace{\medskipamount}}}
\qftype{definition}{Definition}\qftype{convention}{Convention}\qftype{result}{Result}
\qftype{model}{Model}\qftype{method}{Method}\qftype{mechanism}{Mechanism}
\qftype{evidence}{Evidence}\qftype{data}{Data}\qftype{rule}{Rule}
\qftype{pitfall}{Pitfall}\qftype{example}{Example}
\newcommand{\qffield}[1]{\par\textbf{#1.}\enspace}   % starts one of an entry's fields
% \citep[locator]{key} links the key to its line in the sources if it has one,
% and otherwise to its entry in the pinned catalogue PDF.
\newcommand{\qfcatalogue}{../guide/institutional_quantitative_finance_reading_list.pdf}
\NewDocumentCommand{\citep}{o o m}{[\IfValueT{#2}{#1 }%
  \ifcsname r@src:#3\endcsname\hyperref[src:#3]{\nolinkurl{#3}}%
  \else\href{\qfcatalogue\##3}{\nolinkurl{#3}}\fi
  \IfValueTF{#2}{, #2}{\IfValueT{#1}{, #1}}]}
% \qfsource{key}{reference} records one source outside the catalogue, in front/sources.tex.
\newcommand{\qfsource}[2]{\par\noindent\phantomsection\label{src:#1}\nolinkurl{#1}\quad #2\par}
% \qfsections{key, ...} inputs <key>/section.tex for each listed key that has one, in order.
\ExplSyntaxOn
\NewDocumentCommand{\qfsections}{m}
  {\clist_map_inline:nn {#1} {\InputIfFileExists{##1/section}{}{}}}
\ExplSyntaxOff

\newcommand{\figdir}{../../../compendium/execution-market-making/figures}
\pagestyle{plain}
\begin{document}
\chapter*{Market making: legacy derivations}

\section*{Grossman--Miller model, with trading fees and liquidity measures}
% From compendium/legacy/execution-market-making/05_market_making_and_liquidity_provision/01_economics_of_market_making.tex, lines 67--222 (before migration).

\begin{model} \hlt{Grossman-Miller Model} \citep{grossman-miller-1988}\\
Grossman-Miller formalises market making as the supply of \hlt{immediacy} when natural buyers and sellers do not arrive at the same time. Market maker does not want inventory; it accepts one side now because it expects the opposite side later. The price concession is the insurance premium paid by impatient liquidity traders for transferring interim price risk to market makers.

\par\smallskip\noindent\hlt{Assumptions}\\
\begin{enumerate}[label=\roman*.]
\setlength{\itemsep}{0pt}
\item \hlt{timing}: there are three dates $t\in\{1,2,3\}$. Date $1$ is the first liquidity trade, date $2$ is the arrival of the offsetting liquidity trader, and date $3$ is terminal cash settlement;
\item \hlt{assets}: cash earns zero return and the risky asset has terminal value $S_3=\mu+\epsilon_2+\epsilon_3$;
\item \hlt{shocks}: $\epsilon_2,\epsilon_3\overset{\text{iid}}{\sim}N(0,\sigma^2)$. The shock $\epsilon_{t}$ is revealed between dates $t$ and $t-1$;
\item \hlt{participants}: $n$ identical market makers start with zero risky-asset inventory. Liquidity trader $LT_1$ has initial endowment $i$, while $LT_2$ has initial endowment $-i$;
\item \hlt{order imbalance}: if $i>0$, $LT_1$ is a natural seller at date $1$ and $LT_2$ is later buyer. If $i<0$, signs reverse;
\item \hlt{preferences}: all agents have CARA utility $U(X)=-\exp(-\gamma X)$ with risk-aversion parameter $\gamma>0$;
\item \hlt{frictions}: there are no explicit transaction costs, funding costs, or holding costs. The only cost of supplying liquidity is bearing risky inventory between dates $1$ and $2$.
\end{enumerate}

\begin{center}
\refstepcounter{figure}\label{fig:grossman_miller_trading}
\scalebox{0.82}{\input{\figdir/grossman_miller_trading.tikz}}\\
{\small Figure \thefigure: Trading and price setting in the Grossman-Miller model.}
\end{center}

\par\smallskip\noindent\hlt{Setup}\\
Figure~\ref{fig:grossman_miller_trading} summarises the inventory path and price-setting sequence. Let $P_1$ be the date-$1$ trade price and $S_2=\mu+\epsilon_2$ the date-$2$ efficient price after the first shock is known. If agent $j$ holds $q_t^j$ risky units after trading at date $t$, then the date-$2$ and date-$3$ wealth updates are
\begin{align*}
X_2^j&=X_1^j+q_1^j(S_2-P_1), \qquad X_3^j=X_2^j+q_2^j(S_3-S_2).
\end{align*}
At date $2$, conditional on $\epsilon_2$, CARA-normal utility reduces the problem to a certainty-equivalent maximisation:
\begin{align*}
\mathbb{E}\left[U(X_3^j)\mid\epsilon_2\right]
&=-\exp\left(-\gamma\left(X_2^j+q_2^j(\mathbb{E}[S_3\mid\epsilon_2]-S_2)\right)+\frac{1}{2}\gamma^2(q_2^j)^2\sigma^2\right),\\
q_2^{j,\star}
&=\operatorname*{arg\,max}_{q_2^j}\left\{X_2^j+q_2^j(\mathbb{E}[S_3\mid\epsilon_2]-S_2)-\frac{1}{2}\gamma(q_2^j)^2\sigma^2\right\}\\
&=\operatorname*{arg\,max}_{q_2^j}\left\{q_2^j(\mathbb{E}[S_3\mid\epsilon_2]-S_2)-\frac{1}{2}\gamma(q_2^j)^2\sigma^2\right\}.
\end{align*}
The last line drops $X_2^j$ because it is known at date $2$. The first-order condition gives
\begin{align*}
0&=\mathbb{E}[S_3\mid\epsilon_2]-S_2-\gamma q_2^{j,\star}\sigma^2,\\
q_2^{j,\star}&=\frac{\mathbb{E}[S_3\mid\epsilon_2]-S_2}{\gamma\sigma^2}.
\end{align*}
At date $2$, the $n$ MMs, $LT_1$, and $LT_2$ are all present. Market clearing gives
\begin{align*}
nq_1^{MM}+q_1^{LT1}+q_1^{LT2}
&=nq_2^{MM}+q_2^{LT1}+q_2^{LT2}\\
&=(n+2)\frac{\mathbb{E}[S_3\mid\epsilon_2]-S_2}{\gamma\sigma^2}.
\end{align*}
Since the entering date-$2$ holdings sum to zero, $nq_1^{MM}+q_1^{LT1}+q_1^{LT2}=i-i=0$. Hence
\begin{align*}
S_2&=\mathbb{E}[S_3\mid\epsilon_2]=\mu+\epsilon_2, \qquad q_2^j=0.
\end{align*}
Date $2$ is therefore the natural inventory-unwind date: matching liquidity has arrived and no agent wants residual risky inventory.

\par\smallskip\noindent\hlt{Equilibrium}\\
At date $1$, $LT_2$ is absent. The $n$ MMs and $LT_1$ anticipate the efficient date-$2$ unwind, so the same CARA-normal argument gives the desired holding of each present participant:
\begin{align*}
q_1^{j,\star}&=\frac{\mathbb{E}[S_2]-P_1}{\gamma\sigma^2}=\frac{\mu-P_1}{\gamma\sigma^2}.
\end{align*}
Market clearing among $LT_1$ and the $n$ MMs is
\begin{align*}
nq_0^{MM}+q_0^{LT1}
&=nq_1^{MM}+q_1^{LT1}=(n+1)\frac{\mu-P_1}{\gamma\sigma^2}.
\end{align*}
Using $q_0^{MM}=0$ and $q_0^{LT1}=i$ gives
\begin{align*}
i&=(n+1)\frac{\mu-P_1}{\gamma\sigma^2}, \qquad q_1^{j,\star}=\frac{i}{n+1}, \qquad P_1=\mu-\frac{i}{n+1}\gamma\sigma^2.
\end{align*}
The absolute liquidity premium, measured as the concession from efficient value, is
\begin{align*}
p&=\abs{\mu-P_1}=\frac{\abs{i}}{n+1}\gamma\sigma^2.
\end{align*}
If $i>0$, the market makers buy from $LT_1$ at the discounted bid $P_1^b=\mu-p$. If $i<0$, they sell to $LT_1$ at the premium ask $P_1^a=\mu+p$. The implied inside spread is
\begin{align*}
P_1^a-P_1^b&=2p=\frac{2\abs{i}}{n+1}\gamma\sigma^2.
\end{align*}
For $i>0$, $LT_1$ sells $\frac{n}{n+1}i$ units immediately and keeps $\frac{i}{n+1}$ units until date $2$; each market maker holds $\frac{i}{n+1}$ units. The date-$2$ arrival of $LT_2$ clears the remaining inventory at $S_2=\mu+\epsilon_2$.

\par\smallskip\noindent\hlt{Interpretation}\\
The price concession is not an adverse-selection charge. There is no informed trader in the simple model; the order imbalance is temporary and total desired trade nets to zero across dates.\\
The concession is instead a risk-transfer price: the impatient trader replaces an uncertain future execution price with a certain immediate price, while market makers accept exposure to $\epsilon_2$ until matching liquidity arrives.

\par\smallskip\noindent\hlt{Implications}
\begin{enumerate}[label=\roman*.]
\setlength{\itemsep}{0pt}
\item \hlt{larger imbalance}: higher $\abs{i}$ requires a larger concession because more inventory must be warehoused;
\item \hlt{higher volatility}: higher $\sigma^2$ raises the premium because the interim inventory position is riskier;
\item \hlt{higher risk aversion}: higher $\gamma$ raises the premium due to requiring more compensation per unit of risk;
\item \hlt{more market makers}: higher $n$ lowers the premium as imbalance is shared across more risk-bearing capital;
\item \hlt{free entry}: if market presence has fixed opportunity cost $c$, the long-run number of market makers adjusts until risk-adjusted spread capture just covers $c$;
\item \hlt{transitory impact}: a sell imbalance depresses $P_1$ below $\mu$ and then reverses toward $S_2$ when matching demand arrives. Inventory-driven price pressure therefore creates negative short-run return autocorrelation;
\item \hlt{liquidity measure}: the cost is the price of trading now versus waiting, not only the posted bid-ask spread;
\item \hlt{market design}: lower entry, presence, hedging costs increase supply of immediacy. If market-making capacity is scarce, large orders migrate from dealer immediacy to search, block crossing, or delayed execution.
\end{enumerate}

\par\smallskip\noindent\hlt{Trading Costs}\\
Grossman-Miller distinguishes two cost channels. A \hlt{participation cost} $c$ is fixed cost of being present as an active MM, so higher $c$ lowers equilibrium competition $n$ and raises the liquidity premium indirectly. A \hlt{per-share trading fee} $\eta$ is paid when shares trade, so it enters prices directly. At date $2$, fees cancel because remaining positions are liquidated either now or at date $3$:
\begin{align*}
q_2^j
&=\frac{\mathbb{E}[S_3-\eta\mid\epsilon_2]-(S_2-\eta)}{\gamma\sigma^2},
\qquad j\in\{MM,LT1\},\\
q_2^{LT2}
&=\frac{\mathbb{E}[S_3+\eta\mid\epsilon_2]-(S_2+\eta)}{\gamma\sigma^2}.
\end{align*}
Hence $S_2=\mathbb{E}[S_3\mid\epsilon_2]=\mu+\epsilon_2$ as before. At date $1$, for $i>0$, MMs buy from $LT_1$ and expect to sell later; their demand is reduced by the round-trip fee:
\begin{align*}
q_1^{LT1}
&=\frac{\mathbb{E}[S_2-\eta]-(P_1-\eta)}{\gamma\sigma^2}
=\frac{\mu-P_1}{\gamma\sigma^2},\\
q_1^{MM}
&=\frac{\mathbb{E}[S_2-\eta]-(P_1+\eta)}{\gamma\sigma^2}
=\frac{\mu-P_1-2\eta}{\gamma\sigma^2}.
\end{align*}
Market clearing becomes
\begin{align*}
i&=nq_1^{MM}+q_1^{LT1}
=\frac{\mu-P_1}{\gamma\sigma^2}
+n\frac{\mu-P_1-2\eta}{\gamma\sigma^2},
\end{align*}
so
\begin{align*}
P_1&=\mu-\gamma\sigma^2\frac{i}{n+1}-2\left(\frac{n}{n+1}\right)\eta,\\
q_1^{LT1,\star}
&=\frac{i}{n+1}+2\left(\frac{n}{n+1}\right)\frac{\eta}{\gamma\sigma^2}.
\end{align*}
Thus trading fees add an extra liquidity discount and reduce immediacy.\\
$LT_1$ keeps more inventory after date $1$ than in the no-fee case. The initiating liquidity trader pays her own fee directly and pays much of the MMs' round-trip fee indirectly through a lower sale price. If $\abs{i}$ is small relative to $\eta$, the desired trade can fall into a no-trade region.

\par\smallskip\noindent\hlt{Measuring Liquidity}\\
The model can also be read as a reduced-form price-impact model. In an electronic book, $P_1$ is the average execution price of $LT_1$'s marketable sell order against posted liquidity; the initial efficient midpoint is $\mu=\mathbb{E}[S_2]$. Let the date-$1$ quantity sold by $LT_1$ be
\begin{align*}
Q_1^{LT1}&=i-q_1^{LT1,\star}=\frac{n}{n+1}i.
\end{align*}
Then the equilibrium date-$1$ price can be written as
\begin{align*}
P_1&=\mu+\lambda Q_1^{LT1}, \qquad
\lambda=-\frac{\gamma\sigma^2}{n}.
\end{align*}
Here $\lambda$ is the price reaction to signed liquidity demand. A more liquid market has lower $\abs{\lambda}$, which comes from more MM competition $n$, lower risk aversion $\gamma$, or lower interim volatility $\sigma^2$.\\
A second liquidity measure is the autocovariance of price changes. Introduce a date $0$ and public news $\epsilon_1$ before $LT_1$ trades, with efficient price process
\begin{align*}
\mu_0&=\mathbb{E}[S_3], \qquad
\mu_1=\mathbb{E}[S_3\mid\epsilon_1], \qquad
\mu_2=\mathbb{E}[S_3\mid\epsilon_1,\epsilon_2], \qquad
\mu_3=S_3.
\end{align*}
With $S_0=\mu_0$, $S_1=\mu_1+\lambda Q_1^{LT1}$, and $S_2=\mu_2$, price changes are
\begin{align*}
\Delta_1&=S_1-S_0=\epsilon_1+\lambda Q_1^{LT1},\\
\Delta_2&=S_2-S_1=\epsilon_2-\lambda Q_1^{LT1}.
\end{align*}
If liquidity demand is independent of news, then
\begin{align*}
\operatorname{Cov}[\Delta_1,\Delta_2]
&=\operatorname{Cov}[\epsilon_1+\lambda Q_1^{LT1},\epsilon_2-\lambda Q_1^{LT1}]
=-\lambda^2\operatorname{Var}[Q_1^{LT1}]<0.
\end{align*}
Negative price-change autocovariance captures temporary liquidity impact: the first price move contains order pressure, and the next move partially reverses it as price returns toward efficient value. As $\abs{\lambda}\to0$, the observed price process converges to the efficient-price martingale.
\end{model}

\section*{Kyle model}
% From same file, lines 224--336 (before migration).

\begin{model} \hlt{Kyle Model} \citep{s6-kyle}\\
Kyle's static insider-trading model studies how an informed trader monetises a private signal when competitive market makers observe only aggregate order flow. The model is the canonical bridge from \hlt{informational advantage} to \hlt{linear price impact}: informed trades are profitable because liquidity trades provide camouflage, but those same trades move price because market makers learn from order flow.

\par\smallskip\noindent\hlt{Assumptions}\\
\begin{enumerate}[label=\roman*.]
\setlength{\itemsep}{0pt}
\item \hlt{asset value}: one risky asset has terminal value $v\sim N(\mu,\sigma_v^2)$;
\item \hlt{informed trader}: one risk-neutral insider observes $v$ before trading and chooses signed order size $x(v)$;
\item \hlt{liquidity demand}: uninformed liquidity traders submit aggregate order $u\sim N(0,\sigma_u^2)$, independent of $v$ and insensitive to price;
\item \hlt{market makers}: many risk-neutral MMs observe only net order flow $y=x(v)+u$, not the separate components $x(v)$ and $u$;
\item \hlt{competition}: MMs compete order flow profits away, so the transaction price equals conditional expected value, $S(y)=\mathbb{E}[v\mid y]$;
\item \hlt{static auction}: all trading happens once, after which the asset value is realised;
\item \hlt{equilibrium}: the insider maximises expected profit given the price rule, and MMs set prices by Bayes rule given the insider's strategy.
\end{enumerate}

\par\smallskip\noindent\hlt{Setup}\\
The timing is:
\begin{enumerate}[label=\roman*.]
\setlength{\itemsep}{0pt}
\item nature draws $v$ and liquidity demand $u$;
\item the insider observes $v$ and submits $x(v)$;
\item MMs observe only $y=x(v)+u$;
\item MMs set the clearing price $S(y)$;
\item the terminal value $v$ is realised.
\end{enumerate}
The insider's profit is
\begin{align*}
\Pi^{IT}(v,u)&=x(v)\brk{v-S(x(v)+u)}.
\end{align*}
Equilibrium is a fixed point between the insider's trading function and the MMs' pricing function:
\begin{align*}
x(v)&\in\operatorname*{arg\,max}_{x}\mathbb{E}_u\brk{x\brk{v-S(x+u)}\mid v},\\
S(y)&=\mathbb{E}[v\mid y].
\end{align*}
To solve the normal-static version, conjecture a linear equilibrium:
\begin{align*}
S(y)&=\mu+\lambda y,\\
x(v)&=\beta(v-\mu),
\end{align*}
where $\lambda$ is the market's price-impact coefficient and $\beta$ is the insider's trading intensity.

\par\smallskip\noindent\hlt{Equilibrium}\\
Given the pricing rule $S(x+u)=\mu+\lambda(x+u)$, the insider chooses $x$ to maximise conditional expected profit:
\begin{align*}
\max_x\mathbb{E}_u\brk{x\brk{v-\mu-\lambda(x+u)}\mid v} \quad
&=\max_x\brc{x(v-\mu)-\lambda x^2},
\end{align*}
because $\mathbb{E}[u]=0$. The first-order condition gives
\begin{align*}
v-\mu-2\lambda x&=0 \Rightarrow x^\star(v)=\frac{v-\mu}{2\lambda}
\end{align*}
Hence
\begin{align*}
\beta&=\frac{1}{2\lambda}.
\end{align*}
Now impose the MM zero-profit pricing condition. Under $x(v)=\beta(v-\mu)$,
\begin{align*}
y&=\beta(v-\mu)+u.
\end{align*}
Since $(v,y)$ are jointly normal,
\begin{align*}
\mathbb{E}[v\mid y]
&=\mu+\frac{\operatorname{Cov}(v,y)}{\operatorname{Var}(y)}y=\mu+\frac{\beta\sigma_v^2}{\beta^2\sigma_v^2+\sigma_u^2}y.
\end{align*}
Therefore
\begin{align*}
\lambda&=\frac{\beta\sigma_v^2}{\beta^2\sigma_v^2+\sigma_u^2}.
\end{align*}
Combining $\beta=(2\lambda)^{-1}$ with the pricing equation yields
\begin{align*}
\beta^\star&=\frac{\sigma_u}{\sigma_v},\\
\lambda^\star&=\frac{\sigma_v}{2\sigma_u},\\
x^\star(v)&=\frac{\sigma_u}{\sigma_v}(v-\mu),\\
S^\star(y)&=\mu+\frac{\sigma_v}{2\sigma_u}y.
\end{align*}
The insider's conditional expected profit is
\begin{align*}
\mathbb{E}_u[\Pi^{IT}\mid v]
&=x^\star(v)(v-\mu)-\lambda^\star (x^\star(v))^2
=\frac{(v-\mu)^2}{4\lambda^\star},
\end{align*}
and ex-ante expected profit is
\begin{align*}
\mathbb{E}[\Pi^{IT}]&=\frac{\sigma_v\sigma_u}{2}.
\end{align*}
MMs earn zero expected profit conditional on the information they observe:
\begin{align*}
\mathbb{E}\brk{v-S^\star(y)\mid y}=0.
\end{align*}

\par\smallskip\noindent\hlt{Interpretation}\\
The insider does not trade the full mispricing immediately. A larger order extracts more value from information, but it also reveals more information through $y$ and worsens the execution price. The equilibrium order is therefore only a fraction of the private-value gap $v-\mu$.\\
Noise flow is valuable because it makes informed trading hard to identify. Higher $\sigma_u$ lets the insider trade more aggressively and lowers $\lambda^\star$, because a given order-flow imbalance is less informative. Higher $\sigma_v$ raises $\lambda^\star$, because order flow is more value-relevant when private information is more important.\\
The market is \hlt{semi-strong efficient} with respect to MM information: price fully reflects public information plus observed order flow. It is not strong-form efficient because the price does not fully reveal $v$. In equilibrium,
\begin{align*}
\operatorname{Var}(v\mid y)&=\sigma_v^2-\frac{\operatorname{Cov}(v,y)^2}{\operatorname{Var}(y)}
=\frac{1}{2}\sigma_v^2,
\end{align*}
so one static auction reveals only part of the insider's information.

\par\smallskip\noindent\hlt{Implications}
\begin{enumerate}[label=\roman*.]
\setlength{\itemsep}{0pt}
\item \hlt{price impact is informational}: $\lambda^\star$ is not compensation for inventory risk; it is the adverse-selection slope needed when order flow predicts value;
\item \hlt{liquidity traders pay for camouflage}: uninformed flow hides the insider, but loses on average because the insider's profit is financed by liquidity traders while MMs break even;
\item \hlt{more noise lowers impact}: as $\sigma_u$ rises, $\lambda^\star$ falls and larger informed trades can be absorbed without the same price response;
\item \hlt{more private-value risk raises impact}: as $\sigma_v$ rises, order flow becomes more informative and MMs mark prices more aggressively;
\item \hlt{strategic execution matters}: informed traders split the difference between exploiting the signal and concealing it; overly aggressive trading destroys the informational rent through impact;
\item \hlt{permanent impact}: unlike Grossman-Miller inventory pressure, Kyle impact is permanent because price moves toward expected terminal value;
\item \hlt{empirical reading}: a regression of price changes on signed order flow estimates an information-impact coefficient only after separating it from inventory, spread bounce, and mechanical liquidity effects;
\item \hlt{market design}: transparent flow, low noise volume, and concentrated informed trading increase adverse-selection costs; deep uninformed participation makes the market safer for immediacy suppliers.
\end{enumerate}
\end{model}

\section*{Market making using limit orders}
% From same file, lines 338--385 (before migration).

\begin{model} \hlt{Market Making using Limit Orders} \citep{ho-stoll-1981}\\
Grossman-Miller treats immediacy with a single clearing price. In a LOB, the same economics appears through posted limit orders: the MM chooses how far from the midpoint to quote, trading off higher margin per fill against lower execution probability. Figure~\ref{fig:limit_order_market_making} shows the depth and fill-probability objects.

\begin{center}
\refstepcounter{figure}\label{fig:limit_order_market_making}
\scalebox{0.86}{\input{\figdir/limit_order_market_making.tikz}}\\
{\small Figure \thefigure: Limit-order placement and probability of execution.}
\end{center}

\par\smallskip\noindent\hlt{Assumptions}\\
\begin{enumerate}[label=\roman*.]
\setlength{\itemsep}{0pt}
\item \hlt{midpoint}: all traders agree on current value $S_t$;
\item \hlt{orders}: the MM posts sell LO at $S_t+\delta^+$ and buy LO at $S_t-\delta^-$;
\item \hlt{arrivals}: buy and sell MOs arrive with probabilities $p^+$ and $p^-$;
\item \hlt{fill probability}: conditional on MO arrival, execution prob falls with quote distance, $P^\pm(\delta^\pm)=e^{-\kappa^\pm\delta^\pm}$;
\item \hlt{frictions}: no adverse selection, inventory cost, trading cost, or strategic response from other LOs.
\end{enumerate}

\par\smallskip\noindent\hlt{Setup}\\
The MM earns $\delta^+$ when her ask is lifted by a buy MO and $\delta^-$ when her bid is hit by a sell MO. Expected profit per posting decision is
\begin{align*}
\mathbb{E}\left[\Pi(\delta^+,\delta^-)\right]
&=p^+P^+(\delta^+)\delta^+ + p^-P^-(\delta^-)\delta^-\\
&=p^+e^{-\kappa^+\delta^+}\delta^+ + p^-e^{-\kappa^-\delta^-}\delta^-.
\end{align*}

\par\smallskip\noindent\hlt{Equilibrium}\\
The static quoting problem separates across the ask and bid:
\begin{align*}
\max_{\delta^+,\delta^-}\mathbb{E}\left[\Pi(\delta^+,\delta^-)\right]
&=\max_{\delta^+,\delta^-}\left\{p^+e^{-\kappa^+\delta^+}\delta^+ + p^-e^{-\kappa^-\delta^-}\delta^-\right\}.
\end{align*}
The first-order condition on each side is
\begin{align*}
p^\pm e^{-\kappa^\pm\delta^\pm}(1-\kappa^\pm\delta^\pm)&=0,
\end{align*}
so the optimal quote depths are
\begin{align*}
\delta^{\pm,\star}&=\frac{1}{\kappa^\pm}.
\end{align*}

\par\smallskip\noindent\hlt{Interpretation}\\
Posting deeper earns more conditional on execution but reduces probability of execution. Under exponential depth, the optimal quote is placed at the mean depth of incoming market-order walk-up. Arrival probabilities $p^\pm$ scale expected profits, but do not change optimal depths in this static risk-neutral version.

\par\smallskip\noindent\hlt{Limitations}\\
The model isolates the price-probability trade-off, but it is deliberately thin: $P^\pm$, $p^\pm$, and other LOs are exogenous; the MM has no inventory penalty; and there is no adverse selection. More realistic quoting models add dynamic inventory, queue position, cancellation, latency, and informed-flow risk.
\end{model}

\section*{Kyle linear price impact and equilibrium (second legacy version)}
% From compendium/legacy/execution-market-making/03_information_price_formation_and_spreads/05_information_based_models.tex, lines 15--23 (before migration).

\begin{model} \hlt{Kyle Linear Price Impact}\\
Kyle's one-period model separates informed flow from noise flow.\\
The informed trader observes terminal value $p_\infty\sim N(P^\star,\sigma_p^2)$; noise traders submit $n_\omega\sim N(0,\sigma_n^2)$. The market maker observes total flow $y=Q(p_\infty)+n_\omega$ and sets a linear price $f_P(y)=\tilde{P}+\lambda y$.\\
Here $Q$ is informed order size, $\tilde{P}$ is the unconditional price, and $\lambda$ is price impact per unit of net flow. Noise flow is camouflage: the market maker reacts to total flow because it cannot observe which part is informed.
\end{model}

\begin{model} \hlt{Kyle Linear Equilibrium}\\
In the linear equilibrium, $Q(p_\infty)=\frac{p_\infty-\tilde{P}}{2\lambda}$, $\tilde{P}=P^\star$, $\lambda=\frac{\sigma_p}{2\sigma_n}$. The coefficient $\lambda$ is a signal-to-noise price: higher information variance $\sigma_p$ raises impact; higher noise-trader variance $\sigma_n$ hides informed flow and lowers impact.
\end{model}

\section*{Avellaneda--Stoikov framework}
% From compendium/legacy/execution-market-making/05_market_making_and_liquidity_provision/02_avellaneda_stoikov_framework.tex, lines 3--15 and 27--29 (before migration).

\begin{model} \hlt{Avellaneda-Stoikov Setup}\\
A market maker chooses bid and ask quotes while holding inventory $q_t$ in a risky asset with midprice $S_t$. The model balances spread capture against inventory risk under exponential utility and stochastic order arrivals.
\end{model}

\begin{model} \hlt{Reservation Price}\\
With risk aversion $\gamma$, volatility $\sigma$, and horizon $T$, the reservation price is commonly written $r_t=S_t-q_t\gamma\sigma^2(T-t)$.\\
Long inventory lowers reservation price and therefore skews quotes downward; short inventory does the opposite.
\end{model}

\begin{model} \hlt{Symmetric Optimal Spread}\\
In basic exponential-arrival version with arrival sensitivity $k$, total spread is $\psi_t^\star=\gamma\sigma^2(T-t)+\frac{2}{\gamma}\log\left(1+\frac{\gamma}{k}\right)$.\\
The first term prices inventory risk; the second prices execution probability and order-arrival elasticity.
\end{model}

\begin{caveat} \hlt{Framework, Not Production Model}\\
The closed form is a useful control template, but real books require queue position, tick size, discrete inventory limits, latency, rebates, hedging instruments, and regime-dependent fill probabilities.
\end{caveat}


\section*{Glosten--Milgrom and dealer quoting (slice 2, 2026-10-15)}
% Kept verbatim when these blocks were migrated on 2026-10-15 (notes/migration.md, slice 2).
% Destinations: mod:glosten-milgrom, res:glosten-milgrom-half-spreads, res:glosten-milgrom-market-collapse,
% res:adverse-selection-spread-component, def:dealer-inventory-risk, meth:dealer-quote-adjustment.
% Source found for most of this text: Harris, Trading and Exchanges (alpha-harris), chapters "Dealers"
% and "Bid/Ask Spreads", pp. 283-321; the Glosten-Milgrom proposition follows Harris's appendix example
% (pp. 320-321), not the paper. Treat everything here as unchecked own working.

% From compendium/legacy/execution-market-making/03_information_price_formation_and_spreads/05_information_based_models.tex (whole file, deleted 2026-10-15 once empty).
\subsection*{Information-Based Models}

\begin{proposition} \hlt{Glosten-Milgrom Adverse-Selection Spread Component}\\
Let the unconditional value be $V$.\\
The next trader is equally likely to buy or sell; with probability $P$ the trader is informed. If an informed buyer implies value $V+E$ and an informed seller implies value $V-E$, then the zero-cost competitive dealer quotes $a=V+PE$, $b=V-PE$, so the adverse-selection spread component is $a-b=2PE$.
\begin{proof}
Conditional on a buy, the trader is uninformed with probability $1-P$ and informed with probability $P$, so the dealer's conditional value estimate is $(1-P)V+P(V+E)=V+PE$. Conditional on a sell it is $(1-P)V+P(V-E)=V-PE$. Competitive zero-profit quotes equal these conditional values when order-processing and inventory costs are ignored, giving $a-b=2PE$.
\end{proof}
\end{proposition}

\begin{principle} \hlt{Information and Accounting Views Agree}\\
In the simple Glosten-Milgrom setup, $2PE$ is both the conditional-value gap between buyer and seller flow and the expected loss that uninformed traders must pay to compensate the dealer for trading with informed flow.
\end{principle}

% migrated: Kyle Linear Price Impact; Kyle Linear Equilibrium -> mod:kyle-single-auction (full text in notes/artefacts/market-making-lob/derivations.tex)

\begin{principle} \hlt{Adverse Selection Through Flow}\\
The market maker does not need to identify the informed trader. It protects itself by moving price with observed net flow, because aggressive flow is statistically informative.
\end{principle}

% From compendium/legacy/execution-market-making/05_market_making_and_liquidity_provision/02_avellaneda_stoikov_framework.tex (whole file, deleted 2026-10-15 once empty).
\subsection*{Avellaneda-Stoikov Framework}

% migrated: Avellaneda-Stoikov Setup; Reservation Price; Symmetric Optimal Spread -> mod:avellaneda-stoikov (full text in notes/artefacts/market-making-lob/derivations.tex)

% assign: market-making-lob
\begin{taxonomy} \hlt{Quote-Control Segmentation}
\begin{enumerate}[label=\roman*.]
\setlength{\itemsep}{0pt}
\item \hlt{level}: move both quotes with fair value;
\item \hlt{spread}: widen when volatility, toxicity, or uncertainty rises;
\item \hlt{skew}: shift quotes to reduce inventory;
\item \hlt{size}: reduce displayed quantity when adverse selection or limits bind.
\end{enumerate}
\end{taxonomy}

% migrated: Framework, Not Production Model -> mod:avellaneda-stoikov (Pitfalls field) (full text in notes/artefacts/market-making-lob/derivations.tex)

% From compendium/legacy/execution-market-making/05_market_making_and_liquidity_provision/03_adverse_selection_and_toxic_flow.tex (whole file, deleted 2026-10-15 once empty).
\subsection*{Adverse Selection and Toxic Flow}

\begin{definition} \hlt{Adverse Selection Risk}\\
Adverse selection risk is the loss from supplying liquidity to a better-informed trader just before prices move toward that trader's information.
\end{definition}

\begin{mechanism} \hlt{Information Creates Inventory Imbalance}\\
When informed traders buy, dealers sell and become short before prices rise. When informed traders sell, dealers buy and become long before prices fall. The loss is both informational and inventory-driven.
\end{mechanism}

\begin{principle} \hlt{Adverse Selection Spread Component}\\
Dealers set asks conditional on the next trader possibly being an informed buyer, and bids conditional on the next trader possibly being an informed seller.\\
This component rises with information probability, materiality, and trade size.
\end{principle}

\begin{table}[htbp]
\centering
\scriptsize
\stepcounter{table}
\renewcommand{\arraystretch}{1.13}
\setlength{\tabcolsep}{4pt}
\begin{tabularx}{\textwidth}{@{}>{\RaggedRight\arraybackslash}p{0.24\textwidth} >{\RaggedRight\arraybackslash}p{0.31\textwidth} >{\RaggedRight\arraybackslash}X@{}}
\hline
\textbf{Suspected condition} & \textbf{Tactic} & \textbf{Purpose} \\
\hline
Sold to a & Raise ask price; lower ask size & Discourage further sales to informed traders \\
well-informed buyer & Raise bid price; raise bid size & Restore inventory before prices rise \\
& Buy from another trader at the ask & Restore inventory quickly; paying liquidity may cost less than remaining short \\
& Buy a correlated instrument & Hedge inventory risk and speculate on information \\
\hline
Bought from a & Lower bid price; lower bid size & Discourage further purchases from informed traders \\
well-informed seller & Lower ask price; raise ask size & Restore inventory before prices fall \\
& Sell to another trader at the bid & Restore inventory quickly; paying liquidity may cost less than remaining long \\
& Sell a correlated instrument & Hedge inventory risk and speculate on information \\
\hline
\end{tabularx}

\vspace{0.25em}
{\footnotesize Table \thetable. Dealer tactics after trading with a suspected well-informed trader.}
\end{table}

\begin{table}[htbp]
\centering
\scriptsize
\stepcounter{table}
\renewcommand{\arraystretch}{1.13}
\setlength{\tabcolsep}{4pt}
\begin{tabularx}{\textwidth}{@{}>{\RaggedRight\arraybackslash}p{0.15\textwidth} >{\RaggedRight\arraybackslash}p{0.27\textwidth} >{\RaggedRight\arraybackslash}p{0.27\textwidth} >{\RaggedRight\arraybackslash}X@{}}
\hline
\textbf{Buyer is} & \textbf{Raise quotes} & \textbf{Do nothing} & \textbf{Initiate buy at ask} \\
\hline
Well informed & Correct tactic: prevents further informed trading; raises chance of covering before prices rise; reduces realised spread & Mistake: more informed buyers may follow; dealer loses when prices rise; realised spread falls sharply & May be correct if prices soon rise enough; limits later losses, but pays liquidity and reduces realised spread \\
\hline
Uninformed & Mistake: value traders may sell; realised spread and dealer profit fall & Correct tactic: price is presumably still correct; the next seller can pay the full quoted spread & Big mistake: dealer pays unnecessary liquidity and sharply reduces realised spread \\
\hline
\end{tabularx}

\vspace{0.25em}
{\footnotesize Table \thetable. Dealer response consequences after selling to well-informed and uninformed buyers.}
\end{table}

\begin{caveat} \hlt{Informed-Trader Identification}\\
Large, impatient, active, or successful traders may be informed, but these signals are noisy. Dealers therefore price the possibility of information even when no trader can be identified.
\end{caveat}

\begin{caveat} \hlt{Refusing Flow}\\
Dealers can control adverse selection by limiting access, refusing anonymous size, or delaying execution. This protects accepted flow but reduces market-wide immediacy.
\end{caveat}

% From compendium/legacy/execution-market-making/05_market_making_and_liquidity_provision/07_risk_management_for_market_makers.tex (whole file, deleted 2026-10-15 once empty).
\subsection*{Risk Management for Market Makers}

\begin{taxonomy} \hlt{Dealer Inventory Risk Segmentation}
\begin{enumerate}[label=\roman*.]
\setlength{\itemsep}{0pt}
\item \hlt{diversifiable inventory risk}: price changes unrelated to dealer inventory imbalance;
\item \hlt{adverse-selection risk}: price changes inversely related to dealer inventory because informed traders selected the dealer's quote.
\end{enumerate}
The first is reduced by diversification and hedging.\\
The second requires better quoting, client selection, and information inference.
\end{taxonomy}

\begin{principle} \hlt{Dealer Valuation Priority}\\
Dealers care first about market value: the price that attracts balanced flow now. Fundamental value matters because informed traders and value traders force market value toward it.
\end{principle}

\begin{caveat} \hlt{Poorly Informed Dealers}\\
Poorly informed dealers overreact to uninformed flow and underreact to informed flow. They turn inventory quickly, earn smaller realised spreads, and are most vulnerable when client flow looks informative but is not.
\end{caveat}

\begin{mechanism} \hlt{Dealers and Value Traders}\\
Value traders discipline dealer mistakes by trading when quotes move away from fundamental value. Dealers with enough capital and information can wait; constrained dealers must concede price to restore inventory.
\end{mechanism}

% From compendium/legacy/execution-market-making/05_market_making_and_liquidity_provision/01_economics_of_market_making.tex (three blocks).
\subsection*{Economics of Market Making (excerpt)}

\begin{principle} \hlt{Two-Sided Flow}\\
Dealers prefer balanced buy and sell flow. Balanced flow earns spread while limiting inventory; one-sided flow converts market making into risk warehousing.
\end{principle}

\begin{principle} \hlt{Free Entry Bounds Dealer Spreads}\\
With dealer competition and public limit-order competition, spreads are pulled toward the expected cost of providing immediacy. Monopoly power can widen spreads only when entry, routing, or order-driven competition is weak enough to stop clients from switching liquidity suppliers.
\end{principle}

\begin{taxonomy} \hlt{Dealer Spread Cost Segmentation}
\begin{enumerate}[label=\roman*.]
\setlength{\itemsep}{0pt}
\item \hlt{operating costs}: systems, staff, clearing, settlement, membership, and compliance;
\item \hlt{capital costs}: inventory funding, margin, balance-sheet use, and required return on capital;
\item \hlt{inventory losses}: price moves before offsetting flow or hedging;
\item \hlt{informed-flow losses}: trading against clients whose orders predict value changes.
\end{enumerate}
\end{taxonomy}

% From compendium/legacy/execution-market-making/05_market_making_and_liquidity_provision/04_optimal_quoting_and_inventory_control.tex (definition, table and principle).
\subsection*{Optimal Quoting and Inventory Control (excerpt)}

\begin{definition} \hlt{Inventory Imbalance}\\
Let $q_t$ be dealer inventory and $q_t^\star$ the target inventory. The imbalance is $q_t-q_t^\star$. Dealers quote to earn spread while keeping this imbalance financeable.
\end{definition}

\begin{table}[htbp]
\centering
\scriptsize
\stepcounter{table}
\renewcommand{\arraystretch}{1.13}
\setlength{\tabcolsep}{4pt}
\begin{tabularx}{\textwidth}{@{}>{\RaggedRight\arraybackslash}p{0.27\textwidth} >{\RaggedRight\arraybackslash}p{0.34\textwidth} >{\RaggedRight\arraybackslash}X@{}}
\hline
\textbf{Condition} & \textbf{Tactic} & \textbf{Purpose} \\
\hline
Inventories too low or & Raise bid price; increase bid size & Encourage clients to sell \\
clients are net buyers & Raise ask price; decrease ask size & Discourage clients from buying \\
& Take another trader's offer & Immediately raise inventories \\
& Buy a correlated instrument & Hedge inventory risk \\
\hline
Inventories too high or & Lower ask price; increase ask size & Encourage clients to buy \\
clients are net sellers & Lower bid price; decrease bid size & Discourage clients from selling \\
& Hit another trader's bid & Immediately lower inventories \\
& Sell a correlated instrument & Hedge inventory risk \\
\hline
\end{tabularx}

\vspace{0.25em}
{\footnotesize Table \thetable. Dealer tactics for managing inventories and order flow.}
\end{table}

\begin{principle} \hlt{Inventory Control Is Price Discovery}\\
A dealer's search for quotes that balance buy and sell flow also reveals market value. Fundamental value is often learned as a by-product of finding prices at which two-sided flow returns.
\end{principle}


\section*{Spread decomposition and inventory models (slice 3, 2026-10-15)}
% Kept verbatim when these blocks were migrated on 2026-10-15 (notes/migration.md, slice 3).
% Destinations: res:adverse-selection-spread-component, def:dealer-inventory-risk,
% mod:inventory-based-quotes, def:spread-measures, res:roll-spread-estimator,
% res:spread-component-identification, res:spread-determinants, def:outside-spread.
% Sources found for this text: Harris, Trading and Exchanges (alpha-harris), chapters "Dealers",
% "Bid/Ask Spreads", "Value Traders", "Volatility" and "Liquidity and Transaction Cost Measurement";
% Hasbrouck, Empirical Market Microstructure (s6-hasbrouck), chapters 3, 4, 8, 11 and 14. Two claims
% here were found wrong or unsupported (the "twice the mark-to-market loss" proposition and the
% direction of the informed-trading bias in the Roll caveat); see the entries' Pitfalls.
% Treat everything here as unchecked own working.

% From compendium/legacy/execution-market-making/03_information_price_formation_and_spreads/03_bid_ask_spread_decomposition.tex (whole file, deleted 2026-10-15 once empty).
\subsection*{Bid-Ask Spread Decomposition}

\begin{definition} \hlt{Spread Notation}\\
Let $b_t$ be the best bid, $a_t$ the best ask, and $m_t=(a_t+b_t)/2$ the quoted midpoint. The quoted spread is $s_t=a_t-b_t$, the half-spread is $h_t=s_t/2$, and the relative spread is $s_t/m_t$.\\
Cross-sectional spread work should usually compare relative spreads, not dollar spreads.
\end{definition}

\begin{taxonomy} \hlt{Spread Measurement Segmentation}
\begin{enumerate}[label=\roman*.]
\setlength{\itemsep}{0pt}
\item \hlt{quoted spread}: posted ask minus posted bid;
\item \hlt{inside spread}: best ask across dealers or venues minus best bid;
\item \hlt{outside spread}: value trader's sell threshold minus buy threshold;
\item \hlt{effective spread}: execution price relative to the contemporaneous midpoint;
\item \hlt{realised spread}: execution price relative to the midpoint after later price adjustment.
\end{enumerate}
Quoted spreads measure offered liquidity.\\
Realised spreads measure what liquidity suppliers keep after prices move.
\end{taxonomy}

\begin{definition} \hlt{Outside Spread}\\
If a value trader estimates $\widehat{V}$, buys only below $b^o$, and sells only above $a^o$, then $s^o=a^o-b^o$ is the \hlt{outside spread}. It is a private willingness interval, not necessarily a displayed quote. It widens with valuation uncertainty, research cost, impact, financing cost, and adverse-selection risk.
\end{definition}

\begin{taxonomy} \hlt{Spread Component Segmentation}
\begin{enumerate}[label=\roman*.]
\setlength{\itemsep}{0pt}
\item \hlt{order-processing}: systems, clearing, capital, and operating costs;
\item \hlt{inventory}: compensation for warehousing price and funding risk;
\item \hlt{adverse selection}: compensation for trading before prices adjust to informed flow.
\end{enumerate}
\end{taxonomy}

\begin{model} \hlt{Simple Dealer Quote Decomposition}\\
Let $V_0$ be the dealer's current value estimate. Let $V^B$ be the value estimate conditional on the next trader being a buyer, and $V^S$ the estimate conditional on the next trader being a seller. If $C$ is the transaction-cost component and $A=V^B-V^S$ is the adverse-selection component, then $a=V^B+\frac{C}{2}$, $b=V^S-\frac{C}{2}$, $a-b=A+C$.\\
The transaction-cost component is the part a dealer expects to keep in uninformed two-sided flow. The adverse-selection component is the price revision forced by the information in trade direction.
\end{model}

\begin{proposition} \hlt{Adverse-Selection Component Equivalence}\\
Suppose $V_0=(V^B+V^S)/2$ and $V^B\geq V^S$. The adverse-selection spread component $A=V^B-V^S$ is both the gap between conditional value estimates and twice the dealer's mark-to-market loss from quoting around $V_0$ against a correctly signed informed trade.
\begin{proof}
The information view gives $A=V^B-V^S$. If the next buyer is informed and the dealer sells around $V_0$, the value revision is $V^B-V_0=A/2$. If the next seller is informed, the symmetric revision is $V_0-V^S=A/2$. Thus the accounting loss per side is one half of the same component.
\end{proof}
\end{proposition}

\begin{principle} \hlt{Quoted vs Realised Spread}\\
The quoted spread is an offer. The realised spread is smaller when prices move against the liquidity supplier after the trade, which is the empirical footprint of inventory pressure or adverse selection.
\end{principle}

\begin{model} \hlt{Effective Spread Decomposition}\\
For trade sign $\epsilon_t\in\{+1,-1\}$, execution price $p_t$, midpoint $m_t$, and later midpoint $m_{t+\Delta}$, $ES_t=2\epsilon_t(p_t-m_t)$, $RS_t=2\epsilon_t(p_t-m_{t+\Delta})$, $PI_t=2\epsilon_t(m_{t+\Delta}-m_t)$. Then $ES_t=RS_t+PI_t$. $RS_t$ estimates what the liquidity supplier keeps after price adjustment; $PI_t$ estimates the signed price revision associated with the trade.
\end{model}

\begin{mechanism} \hlt{Component Identification}\\
Transaction-cost spread creates bid-ask bounce: trade prices reverse between bid and ask when value is unchanged. Adverse selection creates signed midpoint revision: buyer-initiated trades are followed by higher midpoints and seller-initiated trades by lower midpoints.\\
Empirical decompositions therefore require trade signs and future quote revisions, not prints alone.
\end{mechanism}

\begin{proposition} \hlt{Roll Serial-Covariance Spread Estimator}\\
Assume transaction price $P_t=V_t+\frac{S}{2}Q_t$, where $V_t=V_{t-1}+\epsilon_t$, $\mathbb{E}[\epsilon_t]=0$, $\operatorname{Var}(\epsilon_t)=\sigma^2$, and $Q_t\in\{+1,-1\}$ is an independent buyer/seller sign with equal probabilities. Then $\operatorname{Var}(\Delta P_t)=\sigma^2+\frac{S^2}{2}$, $\operatorname{Cov}(\Delta P_t,\Delta P_{t-1})=-\frac{S^2}{4}$.\\
Hence $S=2\sqrt{-\operatorname{Cov}(\Delta P_t,\Delta P_{t-1})}$ when the serial covariance is negative.
\begin{proof}
Since $\Delta P_t=\epsilon_t+\frac{S}{2}(Q_t-Q_{t-1})$, independence gives $\operatorname{Var}(\Delta P_t)=\sigma^2+\frac{S^2}{4}\mathbb{E}[(Q_t-Q_{t-1})^2]=\sigma^2+\frac{S^2}{2}$.\\
The innovation terms have zero cross-covariance, so $\operatorname{Cov}(\Delta P_t,\Delta P_{t-1}) =\frac{S^2}{4}\mathbb{E}[(Q_t-Q_{t-1})(Q_{t-1}-Q_{t-2})] =-\frac{S^2}{4}$.\\
Solving gives the estimator.
\end{proof}
\end{proposition}

\begin{caveat} \hlt{Roll Estimator Limits}\\
Roll's estimator identifies spread-induced transitory volatility only under strong assumptions. Order splitting, inventory control, serially correlated trade signs, price discreteness, hidden liquidity, and informed trading can make the covariance less negative or positive.
\end{caveat}

\begin{principle} \hlt{Uninformed Trading Loss Channels}\\
Uninformed market orders pay for immediacy through the spread. Uninformed limit orders save the spread only when they execute, and execution is most likely when the quote has become attractive to better-informed flow. The order type changes the loss channel; it does not remove adverse selection.
\end{principle}

\begin{taxonomy} \hlt{Primary Spread Determinants}
\begin{enumerate}[label=\roman*.]
\setlength{\itemsep}{0pt}
\item \hlt{asymmetric information}: raises the adverse-selection component;
\item \hlt{volatility}: raises inventory risk, stale-quote risk, and the value of timing options;
\item \hlt{utilitarian trading interest}: lowers spreads when activity lets dealers and limit-order traders spread fixed costs over many trades.
\end{enumerate}
\end{taxonomy}

\begin{taxonomy} \hlt{Outside Spread Determinant Segmentation}
\begin{enumerate}[label=\roman*.]
\setlength{\itemsep}{0pt}
\item \hlt{valuation error}: wider when fair value is hard to estimate or model errors are correlated;
\item \hlt{winner's curse}: wider when execution itself reveals the trader may be too optimistic or pessimistic;
\item \hlt{research recovery}: wider when few analysed instruments become tradable opportunities;
\item \hlt{position and financing}: wider when expected holding period, leverage, borrow, or funding costs are large;
\item \hlt{market impact}: wider when entering or exiting the position moves price.
\end{enumerate}
Outside spreads are usually wider than dealer spreads because value traders hold risk longer and cannot recycle capital through balanced two-sided flow.
\end{taxonomy}

\begin{taxonomy} \hlt{Order-Driven Spread Determinant Segmentation}
\begin{enumerate}[label=\roman*.]
\setlength{\itemsep}{0pt}
\item \hlt{information asymmetry}: large effect, usually widening spreads;
\item \hlt{cancellation speed}: slow repricing widens spreads by increasing stale-quote risk;
\item \hlt{volatility}: widens spreads through both risk and timing-option value;
\item \hlt{limit-order management cost}: widens spreads when passive liquidity is expensive to monitor;
\item \hlt{trader time value}: widens spreads when waiting is costly;
\item \hlt{commission difference}: passes through roughly one-for-one into the spread;
\item \hlt{risk aversion}: widens spreads when passive execution risk is disliked.
\end{enumerate}
\end{taxonomy}

\begin{taxonomy} \hlt{Secondary Spread Proxy Segmentation}
\begin{enumerate}[label=\roman*.]
\setlength{\itemsep}{0pt}
\item \hlt{information proxies}: disclosure quality, analyst coverage, information vendors, firm age, business complexity, insider-trading enforcement, and expected material news;
\item \hlt{volatility proxies}: leverage, operating uncertainty, storage or inventory risk, macro rate risk, credit risk, and growth uncertainty;
\item \hlt{utilitarian-interest proxies}: volume, trade count, firm size, debt issue size, hedging demand, and availability of close risk substitutes.
\end{enumerate}
These are proxies, not primitives.\\
A variable can load on more than one determinant, so cross-sectional tests need controls.
\end{taxonomy}

\begin{caveat} \hlt{Liquidity and Capital Structure}\\
When asymmetric information is severe and natural trading interest is weak, the spread required for public trading can be too wide for a useful market.\\
Private debt, venture capital, relationship lending, and monitored equity can dominate public trading because they replace anonymous liquidity with due diligence and control rights.
\end{caveat}

% From compendium/legacy/execution-market-making/03_information_price_formation_and_spreads/04_inventory_models.tex (whole file, deleted 2026-10-15 once empty).
\subsection*{Inventory Models}

\begin{principle} \hlt{Inventory Spread Component}\\
Inventory risk is compensation for holding unwanted positions while waiting for offsetting flow.\\
It is separate from adverse selection: inventory risk can exist even when no trader is better informed, because prices can move before the dealer unwinds.
\end{principle}

\begin{taxonomy} \hlt{Inventory Cost Segmentation}
\begin{enumerate}[label=\roman*.]
\setlength{\itemsep}{0pt}
\item \hlt{price risk}: mark-to-market loss while inventory is held;
\item \hlt{funding risk}: financing cost of long inventory and borrow or recall risk for shorts;
\item \hlt{capital risk}: balance-sheet, margin, and risk-limit consumption;
\item \hlt{hedging cost}: basis risk and transaction costs from reducing exposure elsewhere.
\end{enumerate}
\end{taxonomy}

% assign: market-making-lob
\begin{model} \hlt{Inventory-Skew Quoting}\\
Let $q_t$ be inventory, $q^\star$ the target inventory, and $r_t=m_t-\kappa(q_t-q^\star)$ the dealer's reservation price.\\
Quotes can be written as $b_t=r_t-\frac{s_t}{2}$, $a_t=r_t+\frac{s_t}{2}$.\\
When $q_t>q^\star$, the dealer lowers the reservation price to attract buyers and discourage sellers; when $q_t<q^\star$, the dealer raises it to attract sellers and discourage buyers.
\end{model}

\begin{caveat} \hlt{Inventory vs Information Identification}\\
Inventory control and adverse selection can produce the same quote movement after a buy or sell.\\
Identification needs additional information: dealer inventory, hedges, client type, future midpoint revision, or whether quotes revert after inventory is restored.
\end{caveat}

\section*{Volatility, liquidity and crashes (slice 4, 2026-10-15)}
% From compendium/legacy/execution-market-making/03_information_price_formation_and_spreads/08_volatility_and_liquidity_interaction.tex,
% lines 4--91 (before migration): the 14 blocks migrated in slice 4, verbatim. The other 8 blocks
% (liquidity spiral, Flash Crash, systemic risk, trade count) remain in that file, assigned to hft-market-design.
% Destinations: res:spread-determinants, def:volatility-components, mech:bubble-crash-dynamics,
% meth:garman-klass-volatility, meth:spread-volatility-regression (see notes/migration.md, slice 4).

\subsection*{Volatility and Liquidity Interaction (excerpt)}

\begin{principle} \hlt{Volatility Is a Liquidity Price}\\
Volatility raises inventory risk, adverse-selection risk, and the probability that displayed liquidity becomes stale. Liquidity providers respond through wider spreads, smaller displayed size, faster cancellations, or withdrawal.
\end{principle}

\begin{taxonomy} \hlt{Volatility Source Segmentation}
\begin{enumerate}[label=\roman*.]
\setlength{\itemsep}{0pt}
\item \hlt{fundamental volatility}: unexpected changes in efficient value;
\item \hlt{transitory volatility}: liquidity demand, bid-ask bounce, inventory pressure, or temporary order imbalances;
\item \hlt{episodic volatility}: isolated jumps around news, auctions, outages, or large orders;
\item \hlt{continuing volatility}: information, inventory, or funding conditions remain unresolved.
\end{enumerate}
Fundamental volatility is necessary for prices to allocate risk and capital.\\
Transitory volatility is a trading-cost symptom.
\end{taxonomy}

\begin{definition} \hlt{Bubbles and Crashes}\\
A bubble is price far above fundamental value; a crash is a rapid large decline. Both are hard to identify in real time because value is latent.
\end{definition}

\begin{taxonomy} \hlt{Bubble-Crash Dynamics Segmentation}
\begin{enumerate}[label=\roman*.]
\setlength{\itemsep}{0pt}
\item \hlt{optimistic repricing}: new technology, policy, or growth stories raise value estimates;
\item \hlt{momentum reinforcement}: trend followers and order anticipators buy because others are buying;
\item \hlt{valuation restraint}: value traders may be unwilling or unable to short enough size;
\item \hlt{confidence break}: bad news, failed expectations, or funding pressure changes beliefs;
\item \hlt{forced flow}: margin calls, stops, hedges, and risk limits create one-sided selling.
\end{enumerate}
The crash is often the liquidation phase of an earlier pricing error, not an independent event.
\end{taxonomy}

\begin{mechanism} \hlt{Price Accelerator}\\
Price rises transfer wealth to long optimists and away from short pessimists. If optimists buy more after gains, or shorts cover to control losses, the induced order flow pushes prices further in the same direction. The same feedback reverses during crashes when leveraged longs must sell.
\end{mechanism}

\begin{taxonomy} \hlt{Fundamental Volatility Driver Segmentation}
\begin{enumerate}[label=\roman*.]
\setlength{\itemsep}{0pt}
\item \hlt{information surprises}: earnings, macro data, policy, credit, supply-demand, or technology shocks;
\item \hlt{storage and perishability}: inventories buffer shocks when storage is cheap, but amplify them when storage is costly or impossible;
\item \hlt{valuation uncertainty}: high-growth, opaque, or model-sensitive assets reprice sharply when beliefs update;
\item \hlt{leverage}: small asset-value changes produce larger equity-value changes when debt is fixed;
\item \hlt{political and regulatory risk}: cash flows, property rights change when rules, taxes, or convertibility change.
\end{enumerate}
Expected carrying-cost drift is not volatility by itself; the unexpected change in the carry path is.
\end{taxonomy}

\begin{mechanism} \hlt{Transitory Volatility and Reversal}\\
Transitory volatility appears when impatient uninformed trading pushes transaction prices away from efficient value. Bid-ask bounce creates adjacent reversals; larger order imbalances can reverse over minutes, days, or longer when dealers, arbitrageurs, or value traders absorb the pressure.
\end{mechanism}

\begin{caveat} \hlt{Crash Attribution}\\
A large decline is not automatically market failure. Look for reversal, book overload, forced selling, cross-market breaks, or moves beyond news.
\end{caveat}

\begin{model} \hlt{Volatility Component Decomposition}\\
Write observed price as $P_t=M_t+u_t$, where $M_t$ is the efficient price and $u_t$ is transitory microstructure noise. Then $\Delta P_t=\Delta M_t+\Delta u_t$. Fundamental volatility is variation in $\Delta M_t$; transitory volatility is variation in $\Delta u_t$ and its covariance with $\Delta M_t$.\\
Negative serial covariance in $\Delta P_t$ is a useful warning sign that transitory reversals are present.
\end{model}

\begin{caveat} \hlt{Regulating Volatility}\\
Regulators can reduce some transitory volatility through transparency, circuit breakers, margin, and market-maker obligations. They cannot remove fundamental volatility without suppressing price discovery.
\end{caveat}

\begin{definition} \hlt{Garman-Klass Intraday Volatility}\\
For a fixed intraday interval, the Garman-Klass estimator uses open, high, low, and close prices.\\
Let $O$, $H$, $L$, and $C$ denote the interval open, high, low, and close:
\begin{equation*}
\sigma_{GK}= \sqrt{\frac{1}{2}\left(\log\frac{H}{L}\right)^2-\left(2\log 2-1\right)\left(\log\frac{C}{O}\right)^2}.
\end{equation*}
It is useful for comparing intraday uncertainty across markets when transaction-level volatility is too noisy.
\end{definition}

\begin{model} \hlt{Spread-Volatility Regression}\\
A compact empirical check regresses spread on volatility: $\psi_i=a_0+a_1\sigma_i+\epsilon_i$.\\
Here $\psi_i$ is spread observation, $\sigma_i$ is realised volatility, $a_0$ is baseline spread, $a_1$ is spread sensitivity to volatility, and $\epsilon_i$ is residual noise. $\hat{a}_1 > 0$ is consistent with liquidity providers charging more when uncertainty rises.
\end{model}

\begin{model} \hlt{Log-Log Spread-Volatility Model}\\
If spread and volatility are skewed, use $\log \psi_i=a_0+a_1\log\sigma_i+\epsilon_i$.\\
Returning to levels gives the multiplicative relation $\psi_i=A_0\sigma_i^{a_1}e^{\epsilon_i}$ with $A_0=e^{a_0}$. If $\epsilon_i\sim N(0,\sigma_\epsilon^2)$, the unbiased level predictor is $\hat{\psi}= \hat{A}_0\sigma^{\hat{a}_1}e^{\hat{\sigma}_\epsilon^2/2}$.\\
Here $A_0$ is the level intercept, $\sigma_\epsilon^2$ is residual variance, and hats denote estimates.
\end{model}

\begin{caveat} \hlt{Spread-Volatility Identification}\\
A fitted relation is descriptive unless tick size, time of day, venue, depth, and trade intensity are controlled. In large-tick assets, volatility may move through depth rather than spread.
\end{caveat}

\section*{Information, market efficiency and price discovery (slice 5, 2026-10-15)}
% Kept verbatim when these blocks were migrated on 2026-10-15 (notes/migration.md, slice 5).
% Destinations: def:fundamental-value, res:informed-trading-aggregation, def:market-efficiency,
% def:inside-information, def:price-formation-process, def:intraday-overnight-returns, def:outside-spread.
% Sources found: Harris, Trading and Exchanges (alpha-harris), chaps. "Informed Traders and Market
% Efficiency" (pp. 222-244), "Value Traders" (pp. 338-346) and "Insider Trading" (pp. 584-598);
% Lehalle and Laruelle, Market Microstructure in Practice (alpha-laruelle), Introduction and chap. 2,
% sec. 2.5.1. The aggregation proof's variance step and the no-trade band are own working, not the
% sources'. Treat everything here as unchecked own working; corrections are in the entries.

% From compendium/legacy/execution-market-making/03_information_price_formation_and_spreads/01_information_and_market_efficiency.tex (whole file, deleted 2026-10-15 once empty).
\subsection*{Information and Market Efficiency}

\begin{definition} \hlt{Fundamental Value and Price Noise}\\
The \hlt{fundamental value} $V_t$ is the value implied by all payoff-relevant information. The \hlt{market value} $P_t$ is the tradable price, and $\epsilon_t=P_t-V_t$ is price noise.
\end{definition}

\begin{taxonomy} \hlt{Information Segmentation}
\begin{enumerate}[label=\roman*.]
\setlength{\itemsep}{0pt}
\item \hlt{public}: broadly available news, filings, trades, quotes, and macro data;
\item \hlt{private}: costly or restricted information held by a subset of traders;
\item \hlt{stale}: information already incorporated into prices.
\end{enumerate}
\end{taxonomy}

\begin{definition} \hlt{Material Nonpublic Information}\\
Material nonpublic information is value-relevant information that is not broadly available to the market. The legal test is jurisdiction-specific; the microstructure test is whether the information changes $V_t$, who knows it, and whether price has already incorporated it.
\end{definition}

\begin{taxonomy} \hlt{Inside-Information Holder Segmentation}
\begin{enumerate}[label=\roman*.]
\setlength{\itemsep}{0pt}
\item \hlt{corporate insider}: officer, director, employee, or controlling holder with firm information;
\item \hlt{temporary insider}: lawyer, banker, auditor, printer, consultant, exchange staff, or government employee receiving confidential information through a duty;
\item \hlt{tipper-tippee chain}: information passes through relationship, benefit, or confidence channels before trade;
\item \hlt{misappropriator}: trader owes the duty to the information source rather than to the traded issuer;
\item \hlt{skilled analyst}: trader reaches a private conclusion from public data and independent research.
\end{enumerate}
The last category can be economically informed without using inside information. For alpha research, source provenance matters as much as forecast accuracy.
\end{taxonomy}

\begin{principle} \hlt{Signal Tradability Requires Permission}\\
A signal is not tradable merely because it predicts returns. Data provenance, confidentiality duties, firm restricted lists, compliance approvals, and execution capacity determine whether a signal may be used.
\end{principle}

\begin{taxonomy} \hlt{Market Efficiency Segmentation}
\begin{enumerate}[label=\roman*.]
\setlength{\itemsep}{0pt}
\item \hlt{weak-form}: prices reflect past prices and volumes;
\item \hlt{semi-strong}: prices reflect public information;
\item \hlt{strong-form}: prices reflect public and private information;
\item \hlt{economic}: prices reflect information that can be acquired and traded on profitably after costs.
\end{enumerate}
\end{taxonomy}

\begin{proposition} \hlt{Forecast Aggregation Reduces Idiosyncratic Error}\\
Let informed trader $i$ estimate $f_i=V+\epsilon_i$, where $\mathbb{E}[\epsilon_i]=0$ and $\epsilon_i$ are independent with common variance $\sigma^2$. If price is the average forecast $P_N=\frac{1}{N}\sum_{i=1}^N f_i$, then $\mathbb{E}[P_N]=V$ and $\operatorname{Var}(P_N-V)=\sigma^2/N$.
\begin{proof}
Linearity gives $\mathbb{E}[P_N]=V$.\\
Since $P_N-V=\frac{1}{N}\sum_i\epsilon_i$ and errors are independent, $\operatorname{Var}(P_N-V)=N\sigma^2/N^2=\sigma^2/N$.
\end{proof}
\end{proposition}

\begin{caveat} \hlt{Correlated Information Error}\\
Aggregation removes independent errors, not shared errors. If informed traders use the same data, model, or stale signal, prices can move confidently in the wrong direction.
\end{caveat}

\begin{caveat} \hlt{Information-Liquidity Trade-Off}\\
Informative prices are not free. Informed trading transfers wealth from slower liquidity traders and market makers, so policies that improve price efficiency can also reduce displayed liquidity.
\end{caveat}

% From compendium/legacy/execution-market-making/03_information_price_formation_and_spreads/02_price_discovery.tex,
% lines 4--56 (before migration): the 9 blocks migrated in slice 5. The other 3 blocks (fragmentation, fast
% liquidity, outage) remain in that file, assigned to market-structure-execution-quality and hft-market-design.
\subsection*{Price Discovery (excerpt)}

\begin{definition} \hlt{Price Formation Process (PFP)}\\
PFP is the route by which order flow, information, and trading rules become transaction prices. Cleaner PFP links prices more closely to fundamentals; noisy PFP lets microstructure frictions dominate short-horizon prices.
\end{definition}

\begin{mechanism} \hlt{Informed Trading Price Discovery}\\
Informed traders buy when $P_t<V_t$ and sell when $P_t>V_t$. Their trades move prices toward their estimates, but profits shrink as prices incorporate the information.
\end{mechanism}

\begin{taxonomy} \hlt{Value-Trader Profit Source Segmentation}
\begin{enumerate}[label=\roman*.]
\setlength{\itemsep}{0pt}
\item \hlt{information change}: new information moves $V_t$, and price adjusts with delay;
\item \hlt{uninformed liquidity demand}: non-value trading moves $P_t$ away from $V_t$;
\item \hlt{dealer inventory pressure}: intermediaries shade quotes to restore target inventory.
\end{enumerate}
The second and third cases make value traders liquidity suppliers.\\
They are paid for standing against temporary price pressure.
\end{taxonomy}

\begin{proposition} \hlt{Value Traders Supply Liquidity Outside the No-Trade Band}\\
Let a value trader estimate fundamental value by $\widehat{V}$ and face all-in half-cost $c>0$. Buying one unit is attractive only when $P<\widehat{V}-c$, and selling one unit is attractive only when $P>\widehat{V}+c$.
\begin{proof}
Expected surplus from buying is $\widehat{V}-P-c$; expected surplus from selling is $P-\widehat{V}-c$. Each trade is attractive exactly when its surplus is positive.
\end{proof}
\end{proposition}

\begin{caveat} \hlt{No-Trade Is an Information Position}\\
Not buying does not imply selling. When $P$ lies inside the value trader's no-trade band, valuation error, market impact, and the winner's curse can make inactivity the optimal informed decision.
\end{caveat}

\begin{principle} \hlt{Microstructure to Macro Prices}\\
Intraday frictions can affect daily risk when they change the closing price, the auction price, or the path used by benchmark investors. The relevant question is not whether intraday noise exists, but whether it survives into the reference prices used outside the trading day.
\end{principle}

\begin{taxonomy} \hlt{Volatility Scale Segmentation}
\begin{enumerate}[label=\roman*.]
\setlength{\itemsep}{0pt}
\item \hlt{intraday volatility}: uncertainty during trading session, (liquidity, spread, and order-flow shocks);
\item \hlt{overnight volatility}: uncertainty between close and next open (information arrival and auction reopening);
\item \hlt{extraday volatility}: close-to-close risk, mixing intraday and overnight components.
\end{enumerate}
\end{taxonomy}

\begin{definition} \hlt{Intraday and Overnight Returns}\\
For day $t$, let $O_t$ be the open price and $C_t$ the close price.\\
Then $r_t^{\text{intra}}=\log(C_t)-\log(O_t)$ and $r_t^{\text{overnight}}=\log(O_t)-\log(C_{t-1})$.\\
Pairwise correlations can then be computed separately on intraday and overnight return series.
\end{definition}

\begin{caveat} \hlt{Correlation Interpretation}\\
Intraday correlations reflect co-movement during continuous price formation.\\
Overnight correlations are closer to auction and news revaluation. Spikes in daily index correlation hence come from different mechanisms at different horizons.
\end{caveat}

\section*{Strategic trading models (slice 6, 2026-10-15)}
% Kept verbatim when these blocks were migrated on 2026-10-15 (notes/migration.md, slice 6).
% Destinations: def:informed-trader-types, res:informed-trading-aggregation, def:market-efficiency,
% def:inside-information, res:insider-trading-deterrence, def:arbitrage, res:arbitrage-risks,
% def:winners-curse, def:order-anticipation, def:bluffing-manipulation.
% Sources found: Harris, Trading and Exchanges (alpha-harris), chaps. 10, 11, 12, 16, 17, 24 and 29.
% The deterrence, no-trade-band and winner's-curse propositions are own working, not the source's.
% Treat everything here as unchecked own working; corrections are in the entries.

% From compendium/legacy/execution-market-making/03_information_price_formation_and_spreads/06_strategic_trading_models.tex (whole file, deleted 2026-10-15 once empty).
\subsection*{Strategic Trading Models}

\begin{taxonomy} \hlt{Informed Trader Style Segmentation}
\begin{enumerate}[label=\roman*.]
\setlength{\itemsep}{0pt}
\item \hlt{value}: estimate total fundamental value from broad information;
\item \hlt{news}: estimate value changes from recent information;
\item \hlt{information-oriented technical}: exploit predictable errors by other traders;
\item \hlt{arbitrage}: trade relative mispricing across instruments with common valuation factors.
\end{enumerate}
\end{taxonomy}

\begin{taxonomy} \hlt{Informed Trading Speed Segmentation}
\begin{enumerate}[label=\roman*.]
\setlength{\itemsep}{0pt}
\item \hlt{aggressive}: trade quickly when information is short-lived or competitors know it;
\item \hlt{stealth}: trade slowly when impact and information leakage dominate urgency;
\item \hlt{inactive}: avoid trading when the price already reflects the information.
\end{enumerate}
\end{taxonomy}

\begin{taxonomy} \hlt{Insider-Trading Profitability Segmentation}
\begin{enumerate}[label=\roman*.]
\setlength{\itemsep}{0pt}
\item \hlt{materiality}: larger value revisions create larger gross opportunity;
\item \hlt{exclusivity}: fewer informed competitors leave more liquidity before price adjustment;
\item \hlt{time to public release}: near-term disclosure makes profits easier to realise but easier to detect;
\item \hlt{liquidity}: deeper markets allow larger positions before impact erodes the signal;
\item \hlt{traceability}: communication trails, relationship links, and unusual timing raise enforcement risk.
\end{enumerate}
The cases with the largest private payoff are often the worst compliance fact pattern: large, exclusive, traceable, and close to a public event.
\end{taxonomy}

\begin{proposition} \hlt{Deterrence Condition for Illegal Information Trading}\\
Let $G$ be expected gross trading gain from material nonpublic information, $K$ execution cost and leakage, $R$ residual price risk, $q$ detection probability, and $F$ penalty or disgorgement.\\
A purely self-interested trader is deterred when $G\le K+R+qF$.
\begin{proof}
Expected private surplus is $G-K-R-qF$. Trading has positive expected surplus only when this quantity is positive; deterrence requires the opposite inequality.
\end{proof}
\end{proposition}

\begin{definition} \hlt{Arbitrage Opportunity}\\
An arbitrage opportunity exists when related instruments trade outside their normal relation after adjusting for carry, conversion, delivery, funding, borrow, and execution costs. The trade is to buy the cheap leg and sell the expensive leg in a hedge portfolio.
\end{definition}

\begin{notation} \hlt{Basis and Arbitrage Spread}\\
Let $B_t$ be the observed basis between a hedge portfolio and its reference leg, and let $B_t^\star$ be the fair basis implied by carry, delivery, conversion, dividends, funding, borrow, and other contract terms.\\
The \hlt{arbitrage spread} is $A_t=B_t-B_t^\star$.
\end{notation}

\begin{proposition} \hlt{No-Trade Band for Basis Arbitrage}\\
Let $K_t\ge0$ be all-in round-trip cost and let $R_t\ge0$ be compensation for residual basis risk. A symmetric basis trade is attractive only when $\abs{A_t}>K_t+R_t$.
\begin{proof}
Trading the convergence of one unit of basis gives expected gross surplus $\abs{A_t}$.\\
Net expected surplus is $\abs{A_t}-K_t-R_t$, which is positive exactly under the stated condition.
\end{proof}
\end{proposition}

\begin{taxonomy} \hlt{Arbitrage Characterisation Segmentation}
\begin{enumerate}[label=\roman*.]
\setlength{\itemsep}{0pt}
\item \hlt{price characterisation}: arbitrageurs enforce relative value and the law of one price;
\item \hlt{quantity characterisation}: arbitrageurs supply liquidity, transform risk, and repackage exposures;
\item \hlt{implementation characterisation}: arbitrageurs earn spread only if execution, funding, hedging is feasible.
\end{enumerate}
The price view explains why mispricing closes.\\
The quantity view explains why the trade is socially useful and why costs bound the spread.
\end{taxonomy}

\begin{taxonomy} \hlt{Arbitrage Type Segmentation}
\begin{enumerate}[label=\roman*.]
\setlength{\itemsep}{0pt}
\item \hlt{pure arbitrage}: basis strongly mean reverting as delivery, conversion, or cash-settlement links instruments;
\item \hlt{speculative arbitrage}: basis expected to mean revert, but hedge portfolio nonstationary or model-dependent;
\item \hlt{risk arbitrage}: payoff depends on a corporate event, legal condition, or deal completion;
\item \hlt{statistical arbitrage}: factor models identify relative mispricing across many instruments.
\end{enumerate}
Pure arbitrage can lose money if convergence is slow, financing fails, or implementation costs exceed spread.
\end{taxonomy}

\begin{taxonomy} \hlt{Pure Arbitrage Mechanism Segmentation}
\begin{enumerate}[label=\roman*.]
\setlength{\itemsep}{0pt}
\item \hlt{shipping}: same asset trades in different locations, bounded by transport and storage costs;
\item \hlt{delivery}: futures and cash prices converge toward delivery or settlement;
\item \hlt{conversion}: one risk form is manufactured from another, such as options, ETFs, or commodity crushes;
\item \hlt{cross-rate}: triangular or multi-leg exchange rates imply an inconsistent synthetic price.
\end{enumerate}
Closest to mechanical arbitrage, but binding constraint is usually speed, balance sheet, or operational capacity.
\end{taxonomy}

\begin{taxonomy} \hlt{Arbitrage Risk Segmentation}
\begin{enumerate}[label=\roman*.]
\setlength{\itemsep}{0pt}
\item \hlt{implementation}: legs fill at worse prices, or one leg fills while another does not;
\item \hlt{basis}: the spread widens before convergence, forcing mark-to-market losses;
\item \hlt{model}: the assumed fair relation is wrong or misses a contract feature;
\item \hlt{carry}: funding, borrow, dividends, storage, insurance, or spoilage costs differ from expectation;
\item \hlt{staying power}: margin, financing, or risk limits force liquidation before convergence;
\item \hlt{specific-factor}: an idiosyncratic shock breaks the relative-value relation.
\end{enumerate}
Question is not only whether the basis is cheap, but whether the portfolio can survive the path to convergence.
\end{taxonomy}

\begin{principle} \hlt{Common vs Specific Factor Test}\\
Relative-value trades work best when the spread move is caused by uninformed liquidity demand or slow common-factor adjustment. They fail when the move reflects a permanent instrument-specific shock.\\
Residualisation is therefore not cosmetic; it decides whether the spread is an opportunity or new information.
\end{principle}

\begin{principle} \hlt{Orthogonal Information Advantage}\\
The best informed traders are not only accurate; their errors are different from other traders' errors. Orthogonal information is more valuable than crowded information with the same average accuracy.
\end{principle}

\begin{principle} \hlt{Value Traders vs News Traders}\\
Value traders profit when price differs from value because liquidity demand or dealer inventory pressure moved price away from fundamentals.\\
News traders profit when value itself has changed and price has not caught up. They trade against each other when one mistakes stale price pressure for mispricing, or stale news for new information.
\end{principle}

\begin{caveat} \hlt{Informed Trading Paradox}\\
Perfectly informative prices eliminate informed profits; zero informed profits stop prices from becoming informative. Markets persist between these limits because information is costly, temporary, and unevenly distributed.
\end{caveat}

\begin{caveat} \hlt{Skilled Speculation vs Inside Information}\\
Correct forecasts from public filings, alternative data with clean rights, channel checks, or independent modelling can be legitimate informed trading. The same trade becomes a compliance problem when the edge comes from material nonpublic information or breach of confidence. Timing alone cannot distinguish the cases.
\end{caveat}

\begin{mechanism} \hlt{Insider-Trading Market Footprint}\\
Leaked or illegal information may appear as abnormal signed volume, option activity, price drift, widened spreads, quote fading, or dark-pool activity before news. These are surveillance leads, not proof: rumours, public-information analysis, hedging, and unrelated liquidity demand can look similar.
\end{mechanism}

\begin{definition} \hlt{Winner's Curse}\\
The winner's curse is the adverse information revealed by winning a common-value trade or auction.\\
A buyer who wins may be the most optimistic bidder; a seller who wins may be the most pessimistic offeror.
\end{definition}

\begin{proposition} \hlt{Winner's-Curse Bid Adjustment}\\
Let a buyer's estimate be $\widehat{V}$ and let $\omega_N=\mathbb{E}[\widehat{V}-V\mid\text{buyer wins among }N\text{ buyers}]$ be the expected overvaluation conditional on winning. With all-in cost $c$, a zero-profit bid is $b^\star=\widehat{V}-\omega_N-c$.
\begin{proof}
Conditional expected profit from bid $b$ is $\mathbb{E}[V-b-c\mid\text{win}]=\widehat{V}-\omega_N-b-c$.\\
Setting this to zero gives $b^\star=\widehat{V}-\omega_N-c$.
\end{proof}
\end{proposition}

\begin{taxonomy} \hlt{Winner's-Curse Severity Segmentation}
\begin{enumerate}[label=\roman*.]
\setlength{\itemsep}{0pt}
\item \hlt{competitor count}: more rival bidders increase the chance that winning selects an extreme estimate;
\item \hlt{valuation dispersion}: noisier estimates require larger bid shading;
\item \hlt{information overlap}: common data reduce dispersion; independent but noisy private estimates increase selection risk;
\item \hlt{market power}: a trader facing few competitors may bid less aggressively and still win.
\end{enumerate}
\end{taxonomy}

\begin{taxonomy} \hlt{Bluffing and Manipulation Segmentation}
\begin{enumerate}[label=\roman*.]
\setlength{\itemsep}{0pt}
\item \hlt{rumourmonger}: spreads false or selectively framed information to make others trade;
\item \hlt{price manipulator}: trades to make prices, volume, or timing look informative;
\item \hlt{wash trader}: arranges offsetting trades with confederates to manufacture activity.
\end{enumerate}
The bluff works only if other traders mistake the signal for information or urgent liquidity demand.
\end{taxonomy}

\begin{principle} \hlt{Bluffers vs Informed Speculators}\\
Both may trade aggressively and advertise conviction. The difference is informational content: informed speculators trade on value-relevant beliefs, while bluffers create signals they hope others will misread.
\end{principle}

\begin{mechanism} \hlt{Value Traders Discipline Bluffs}\\
Value traders break bluffs by trading against clear mispricing. Where valuation, borrow, or funding is constrained, bluffs can persist and attract momentum flow.
\end{mechanism}

\begin{caveat} \hlt{Manipulation Evidence Problem}\\
Price reversals are not enough to prove manipulation. Prosecutors must usually separate a false signal or wash trade from a bad but legitimate valuation view, which is why small or careful bluffs are hard to punish.
\end{caveat}

\begin{taxonomy} \hlt{Order-Anticipation Segmentation}
\begin{enumerate}[label=\roman*.]
\setlength{\itemsep}{0pt}
\item \hlt{front running}: trade before a known or inferred order;
\item \hlt{sentiment-oriented technical}: predict trades uninformed traders will make;
\item \hlt{quote matching}: step ahead of passive quotes to extract their option value;
\item \hlt{squeezing}: corner scarce supply and force constrained traders to liquidate;
\item \hlt{stop gunning}: push prices to trigger clustered stop orders.
\end{enumerate}
These strategies use information about future order flow rather than fundamental value.
\end{taxonomy}

\begin{taxonomy} \hlt{Specialist Order-Flow Strategy Segmentation}
\begin{enumerate}[label=\roman*.]
\setlength{\itemsep}{0pt}
\item \hlt{flow speculation}: trade from early knowledge of routed orders or imbalance;
\item \hlt{quote matching}: improve by the minimum tick to step ahead of public interest;
\item \hlt{cream skimming}: fill less toxic marketable flow and leave harder flow to the book;
\item \hlt{stopping option}: guarantee a price, then choose execution after observing later order flow.
\end{enumerate}
These strategies monetise privileged sequencing, so audit trails and public-order precedence are central controls.
\end{taxonomy}

\begin{mechanism} \hlt{Front-Running Effect}\\
Front runners accelerate anticipated order impact. They can speed price discovery for informed flow, but usually raise execution costs and discourage large-order exposure.
\end{mechanism}

\begin{mechanism} \hlt{Sentiment-Oriented Technical Trading}\\
These traders forecast predictable non-value flow: tax timing, index changes, hedging, funding needs, or behavioural demand. Profits decay when value traders correct the move or the pattern crowds.
\end{mechanism}

\begin{caveat} \hlt{Quote Matching and Tick Protection}\\
Quote matching extracts value from standing limit orders without supplying durable liquidity. Time precedence and meaningful tick sizes protect passive liquidity by making queue jumping costly.
\end{caveat}

\begin{caveat} \hlt{Squeezes and Stop Gunning}\\
Squeezes and stop gunning profit from forced trades. They usually make prices less informative, raise transaction costs, and reduce liquidity for traders who must close positions.
\end{caveat}

\section*{Liquidity and market impact (slice 7, 2026-10-15)}
% Kept verbatim when these blocks were migrated on 2026-10-15 (notes/migration.md, slice 7).
% Destinations: def:price-impact, mod:propagator, res:no-price-manipulation (compendium/execution-trading-costs/,
% pathway price-impact-execution) and def:bluffing-manipulation.
% Sources found: Webster, Handbook of Price Impact Modeling (wide-webster-price-impact), chaps. 1-2;
% Bouchaud, Bonart, Donier and Gould, Trades, Quotes and Prices (catalogue-0528), chaps. 11-13 and 19;
% Harris (alpha-harris), chap. 12. The round-trip proposition is an accounting identity, not a
% no-manipulation result. Treat everything here as unchecked own working; corrections are in the entries.

% From compendium/legacy/execution-market-making/03_information_price_formation_and_spreads/07_liquidity_and_market_impact.tex (whole file, deleted 2026-10-15 once empty).
\subsection*{Liquidity and Market Impact}

\begin{taxonomy} \hlt{Market Impact Lens Segmentation}
\begin{enumerate}[label=\roman*.]
\setlength{\itemsep}{0pt}
\item \hlt{theoretical}: strategic splitting and information inference;
\item \hlt{reduced-form}: price response to signed volume, participation, volatility, and time;
\item \hlt{practitioner}: cost of urgency, benchmark choice, and venue access.
\end{enumerate}
\end{taxonomy}

\begin{definition} \hlt{Temporary and Permanent Impact}\\
Temporary impact is the price concession needed to obtain immediacy and is expected to decay as liquidity replenishes. Permanent impact is the residual price shift after decay, but it may reflect information in the order rather than mechanical pressure.
\end{definition}

\begin{mechanism} \hlt{Propagator Intuition}\\
A trade has an impulse effect on price, and that effect decays through time. The observed price path is the sum of many signed trade impulses with different decay profiles. This separates immediate pressure from later relaxation: a print can move price now without requiring the full move to be permanent.
\end{mechanism}

\begin{proposition} \hlt{Round-Trip Manipulation Profit Test}\\
Suppose a trader completes a zero-inventory round trip, buying and selling total quantity $Q>0$. Let $\bar{P}_B$ be the average buy price, $\bar{P}_S$ the average sale price, and $c$ the all-in cost per share.\\
The round trip is profitable if and only if $\bar{P}_S-\bar{P}_B>c$.
\begin{proof}
The trader pays $Q\bar{P}_B$ on buys and receives $Q\bar{P}_S$ on sales. Net profit after costs is $Q(\bar{P}_S-\bar{P}_B-c)$, which is positive exactly under the stated condition.
\end{proof}
\end{proposition}

\begin{principle} \hlt{Impact Symmetry Disciplines Bluffing}\\
Liquidity suppliers are vulnerable when buys and sells, large and small orders, or fast and slow trades receive inconsistent price impact. A manipulation-resistant rule makes impact depend on net pressure and risk, so a zero-net sequence cannot mechanically create a higher average sale price than buy price after costs.
\end{principle}

\begin{caveat} \hlt{Permanent Impact Identification}\\
Calling a long-run price move permanent impact is risky unless market moves, sector moves, autocorrelated flow, and the signal that caused the trade have been separated.
\end{caveat}

\section*{Market impact models (slice 8, 2026-10-15)}
% Kept verbatim when these blocks were migrated on 2026-10-15 (notes/migration.md, slice 8).
% Destinations (compendium/execution-trading-costs/, pathway price-impact-execution): def:price-impact,
% mod:propagator, ev:square-root-impact, meth:post-execution-impact-profile.
% Sources found: Bacry, Iuga, Lasnier and Lehalle (wide-bacry-life-cycle, arXiv 2014), sec. 7;
% Bouchaud, Gefen, Potters and Wyart (wide-bouchaud-propagator, arXiv 2003); Almgren and Chriss
% (alpha-almgren, preprint 2000); Bouchaud, Bonart, Donier and Gould (catalogue-0528), chaps. 10, 12, 13.
% The propagator block departs from its sources (timing, noise term, volume weight); corrections are in
% the entries. Treat everything here as unchecked own working.

% From compendium/legacy/execution-market-making/04_order_flow_and_empirical_microstructure/05_market_impact_models.tex (whole file, deleted 2026-10-15 once empty).
\subsection*{Market Impact Models}

\begin{taxonomy} \hlt{Model Family Segmentation}
\begin{enumerate}[label=\roman*.]
\setlength{\itemsep}{0pt}
\item \hlt{Almgren-Chriss}: optimise a schedule by balancing expected impact cost against execution risk;
\item \hlt{linear impact}: use tractable instantaneous or temporary cost proportional to trading rate or signed size;
\item \hlt{propagator}: model each signed trade as a decaying impulse in future prices;
\item \hlt{empirical square-root}: summarise metaorder cost as concave in order size relative to market volume.
\end{enumerate}
\end{taxonomy}

\begin{principle} \hlt{Permanent Impact or Price Anticipation}\\
Long-horizon impact mixes metaorder pressure and information revealed by trading. Remove the information effect before calling the residual permanent impact.
\end{principle}

\begin{taxonomy} \hlt{Permanent-Impact View Segmentation}
\begin{enumerate}[label=\roman*.]
\setlength{\itemsep}{0pt}
\item \hlt{mechanical view}: prices move because buying and selling pressure are imbalanced;
\item \hlt{informational view}: prices move because the trade reveals new expectations about value;
\item \hlt{hybrid view}: executions contain both pressure and information, so identification matters more than labels.
\end{enumerate}
\end{taxonomy}

\begin{model} \hlt{Propagator Midprice Dynamics}\\
In trade time, a standard propagator writes the midprice as $m_k=m_{-\infty}+\sum\limits_{k'<k}G(k-k')\epsilon_{k'}+\eta_k$.\\
Here $m_k$ is the midprice after trade $k$, $\epsilon_k\in\{-1,+1\}$ is trade sign, $G(\cdot)$ is the decaying impact kernel, $\eta_k$ is fair-price noise, and $m_{-\infty}$ is an initial price level.
\end{model}

\begin{mechanism} \hlt{Trade-Sign Memory and Kernel Decay}\\
Trade signs are positively autocorrelated because large orders are split and similar agents trade together.\\
A decaying kernel offsets this persistence: impact must decay enough that predictable order flow does not imply predictable price drift.
\end{mechanism}

\begin{model} \hlt{Price Response Curve}\\
For horizon $h$, define the response curve by $\mathcal{R}(h)=\mathbb{E}[(m_{k+h}-m_k)\epsilon_k]$.\\
Here $\mathcal{R}(h)$ measures the expected signed price move after a trade. It is estimated with the trade-sign autocorrelation $C(\ell)=\mathbb{E}[\epsilon_{k+\ell}\epsilon_k]$ and the kernel increments $\Delta G(\ell)=G(\ell+1)-G(\ell)$.
\end{model}

\begin{taxonomy} \hlt{Propagator Kernel Segmentation}
\begin{enumerate}[label=\roman*.]
\setlength{\itemsep}{0pt}
\item \hlt{non-parametric}: estimate discrete kernel increments directly from response and sign autocorrelation;
\item \hlt{exponential}: impose fast decay with a small number of parameters;
\item \hlt{power law}: allow slow decay and long-memory order flow;
\item \hlt{agent-specific}: estimate separate kernels or intensities by participant class.
\end{enumerate}
\end{taxonomy}

\begin{process} \hlt{Debiasing Long-Horizon Impact}
\begin{enumerate}[label=\roman*.]
\setlength{\itemsep}{0pt}
\item classify trades by whether they are likely information-driven or liquidity-driven;
\item remove market and sector moves with a benchmark or factor model;
\item adjust for autocorrelated metaorder flow after the execution day;
\item measure the remaining post-execution drift over several days.
\end{enumerate}
\end{process}

\begin{model} \hlt{CAPM-Like Post-Execution Decomposition}\\
For event day $D$, a stock's post-execution move can be split into systematic and idiosyncratic parts:
\begin{equation*}
\log(P_d)-\log(P_{d-1})=
\beta\left[\log(I_d)-\log(I_{d-1})\right]+\Delta W_d.
\end{equation*}
Here $P_d$ is the stock close, $I_d$ the reference index, $\beta$ the index exposure, and $\Delta W_d$ the residual stock-specific move. Cumulative $\Delta W_d$ isolates the profile after removing index co-movement. The intuition is attribution: do not charge the execution for a market move that the whole index experienced.
\end{model}

\begin{caveat} \hlt{Autocorrelated Flow Bias}\\
Metaorder flow can remain signed and autocorrelated for days. If later same-direction executions are ignored, their impact may be falsely attributed to the original order.
\end{caveat}

\section*{Price impact estimation (slice 9, 2026-10-15)}
% Kept verbatim when these blocks were migrated on 2026-10-15 (notes/migration.md, slice 9).
% Destinations (compendium/execution-trading-costs/, pathway price-impact-execution): def:price-impact,
% ev:intraday-impact-profile, meth:almgren-impact-estimation, meth:order-flow-impact-regression,
% ev:index-reconstitution-impact.
% Sources found: Lehalle and Laruelle (alpha-laruelle), chap. 3, sec. 3.2; Harris (alpha-harris), boxes
% pp. 435 and 487; Hasbrouck (s6-hasbrouck), sec. 9.7-9.8; Bacry et al. (wide-bacry-life-cycle, arXiv 2014);
% Almgren et al. (wide-almgren-direct-impact, preprint 2005). The regression block omits the size-independent
% permanent term of the general Glosten-Harris model; corrections are in the entries.
% Treat everything here as unchecked own working.

% From compendium/legacy/execution-market-making/04_order_flow_and_empirical_microstructure/04_price_impact_estimation.tex (whole file, deleted 2026-10-15 once empty).
\subsection*{Price Impact Estimation}

\begin{definition} \hlt{Execution-Time Impact}\\
Execution-time impact is the price move observed while the order is being traded, measured from the perspective of the order side and often normalised by spread or volatility.
\end{definition}

\begin{empiricalfinding} \hlt{Impact Has a Hump Shape}\\
Intraday impact usually rises during execution and partly decays afterwards. The remaining shift is the candidate permanent component; the decayed part is temporary impact and liquidity replenishment.
\end{empiricalfinding}

\begin{taxonomy} \hlt{Impact Driver Segmentation}
\begin{enumerate}[label=\roman*.]
\setlength{\itemsep}{0pt}
\item \hlt{stock characteristics}: liquidity, volatility, spread, tick size, and typical depth;
\item \hlt{order characteristics}: size, duration, participation rate, side, and aggressiveness;
\item \hlt{timing}: trend, rebalancing pressure, signal arrival, auctions, and market-wide events;
\item \hlt{venue access}: lit depth, dark liquidity, queue priority, and routing constraints.
\end{enumerate}
\end{taxonomy}

\begin{model} \hlt{Participation-Duration Shape}\\
A compact empirical specification writes market impact as increasing in duration and participation:
\begin{equation*}
\text{Market Impact} \propto \text{Duration}^{\alpha}\times \text{Participation}^{\gamma}.
\end{equation*}
Here $\alpha$ is the duration exponent, $\gamma$ is the participation exponent. The exponents are calibrated from executions; in many datasets the participation exponent is concave, so extra participation hurts, but less than linearly. The model is not structural; it is a compact way to express the practitioner trade-off between urgency and footprint.
\end{model}

\begin{model} \hlt{Order-Flow Impact Regression}\\
A Glosten-Harris style reduced form writes $\Delta P_t=\lambda\,\epsilon_t q_t+\psi\,\Delta(\epsilon_t q_t)+\eta\,\Delta\epsilon_t+u_t$ where $\epsilon_t$ is trade sign and $q_t$ is trade size. The level signed-size term estimates permanent information impact; changes in sign or signed size capture transitory spread and liquidity-supply effects.
\end{model}

\begin{empiricalfinding} \hlt{Index Reconstitution Impact}\\
Index additions and deletions create predictable one-sided flow. Price impact is largest when forced index demand is large relative to stock liquidity; partial reversal after the event suggests temporary demand pressure, while the residual move reflects changed ownership, attention, or information.
\end{empiricalfinding}

\begin{caveat} \hlt{Order-Flow Regression Identification}\\
Signed order flow is endogenous: informed traders choose when and how much to trade. Use controls for market moves, time of day, volatility, spread, and order splitting before interpreting coefficients as causal impact.
\end{caveat}

\begin{caveat} \hlt{Timing Is Not Noise}\\
Two orders with identical size and participation can have different impact because one trades with the market and the other trades against it. Impact estimation therefore has to separate execution mechanics from the signal that caused the trade.
\end{caveat}

\section*{Optimal execution problem (slice 10, 2026-10-15)}
% Kept verbatim when these blocks were migrated on 2026-10-15 (notes/migration.md, slice 10).
% Destinations (compendium/execution-trading-costs/, owned by the section overview execution-trading-costs):
% mod:almgren-chriss-execution, def:block-trade.
% Sources found: Lehalle and Laruelle (alpha-laruelle), app. A, sec. A.7; Almgren and Chriss (alpha-almgren,
% author preprint of December 2000); Harris (alpha-harris), chap. 15. The comparative static on impact
% sensitivity and the convexity claim are corrected in the entry. Treat everything here as unchecked own working.

% From compendium/legacy/execution-trading-costs/01_execution_and_algorithmic_trading/01_optimal_execution_problem.tex (whole file, deleted 2026-10-15 once empty).
\subsection*{Optimal Execution Problem}

\begin{principle} \hlt{Impact-Risk Trade-Off}\\
Optimal execution balances two forces: market impact rewards slower trading, while price risk rewards faster completion. A schedule is a compromise between footprint and exposure.
\end{principle}

\begin{model} \hlt{Discrete Execution Schedule}\\
Plan a buy order of size $v^\star$ over $N$ slices. During slice $n$, trade $v_n$ shares, expected market volume is $V_n$, volatility is $\sigma_n$, and the fair price evolves as $S_{n+1}=S_n+\sigma_n\xi_{n+1}$, $\sum\limits_{n=1}^{N}v_n=v^\star$.\\
Here $\xi$ is a standard Gaussian shock. The schedule controls timing; the market supplies uncertainty.
\end{model}

\begin{model} \hlt{Temporary Impact Cost}\\
A simple execution price for slice $n$ is $\widetilde{S}_n(v_n)=S_n+\kappa\sigma_n\left(\frac{v_n}{V_n}\right)^\gamma$, $W=\sum\limits_{n=1}^{N}v_n\widetilde{S}_n(v_n)$.\\
Here $W$ is total execution notional, and $\kappa$ and $\gamma$ are stock-dependent impact parameters.\\
Participation $v_n/V_n$ is the key lever: trading more of the available volume costs more.
\end{model}

\begin{definition} \hlt{Expectation-Minimising Schedule}\\
Ignoring risk, the minimiser allocates more to high-volume and low-volatility slices: $v_n^{\mathrm{exp}} =\frac{V_n/\sigma_n^{1/\gamma}}{\sum\limits_{\ell=1}^{N}V_\ell/\sigma_\ell^{1/\gamma}}v^\star$.\\
If $V_n$ and $\sigma_n$ are constant, the schedule is uniform: $v_n=v^\star/N$.
\end{definition}

\begin{mechanism} \hlt{Risk-Only Schedule}\\
If impact is ignored and only price-risk variance matters, the optimal action is to trade immediately. Waiting leaves remaining inventory exposed to random price moves.
\end{mechanism}

\begin{model} \hlt{Mean-Variance Execution Objective}\\
Let $x_n$ be the remaining quantity after slice $n$, with $x_0=v^\star$ and $x_N=0$.\\
A compact objective is $\mathcal{C} =v^\star S_0 +\kappa\sum\limits_{n=1}^{N}\sigma_n\frac{v_n^{\gamma+1}}{V_n^\gamma} +\lambda\sum\limits_{n=1}^{N}x_n^2\sigma_n^2$.\\
The middle term is impact cost; the last term is inventory risk. Higher risk aversion $\lambda$ front-loads the schedule, while higher impact sensitivity back-loads it.
\end{model}

\begin{taxonomy} \hlt{Execution Schedule Segmentation}
\begin{enumerate}[label=\roman*.]
\setlength{\itemsep}{0pt}
\item \hlt{risk-neutral}: follows volume and volatility to reduce expected impact;
\item \hlt{risk-only}: finishes immediately to remove exposure;
\item \hlt{mean-variance}: trades along a convex remaining-quantity path;
\item \hlt{tactical overlay}: lets dark pools, SORs, alpha, and limits modify child orders inside the schedule.
\end{enumerate}
\end{taxonomy}

\begin{definition} \hlt{Block Order}\\
A block order is too large to complete easily through normal displayed liquidity.\\
Fixed-share thresholds are venue conventions; for execution research, a block is better measured as a fraction of average daily volume, typical depth, or expected participation capacity.
\end{definition}

\begin{taxonomy} \hlt{Block Execution Problem Segmentation}
\begin{enumerate}[label=\roman*.]
\setlength{\itemsep}{0pt}
\item \hlt{latent demand}: natural contra interest exists but has not posted orders;
\item \hlt{order exposure}: revealing the parent attracts anticipators and liquidity withdrawal;
\item \hlt{price discrimination}: liquidity suppliers wait because they fear selling too cheaply or buying too richly;
\item \hlt{asymmetric information}: suppliers demand concessions when the initiator may be informed.
\end{enumerate}
Block execution buys liquidity with information: disclose enough to find size, but not enough to be exploited.
\end{taxonomy}

\end{document}
